\documentclass[11pt]{article}
\usepackage[a4paper,margin=25mm]{geometry}
\usepackage[T1]{fontenc}
\usepackage[utf8]{inputenc}
\usepackage{lmodern}
\usepackage{microtype}
\usepackage{booktabs,tabularx,array,longtable}
\usepackage{enumitem}
\usepackage{xcolor}
\usepackage{titlesec}
\usepackage{float}
\PassOptionsToPackage{hyphens}{url}
\usepackage[hidelinks,breaklinks]{hyperref}
\urlstyle{same}

\definecolor{ocsi}{RGB}{30,60,90}
\titleformat{\section}{\normalfont\large\bfseries\color{ocsi}}{\thesection}{0.6em}{}
\titleformat{\subsection}{\normalfont\normalsize\bfseries}{\thesubsection}{0.6em}{}
\setlist{nosep,leftmargin=*}
\newcommand{\tphe}{T\mbox{-}PHE}
\newcolumntype{Y}{>{\raggedright\arraybackslash}X}
\renewcommand{\thesection}{S\arabic{section}}
\renewcommand{\thetable}{S\arabic{table}}

\begin{document}

\begin{center}
{\LARGE\bfseries\color{ocsi} Online Supplementary Material}\par\vspace{2mm}
{\large Before Family Risk Becomes a Health Crisis: Trustworthy Premarital Health Education as Primary Prevention Infrastructure}\par\vspace{2mm}
{\normalsize Yaoharee Lahtee$^{1,2}$}\par
{\small $^{1}$ARAYA Nikah Social Enterprise Co., Ltd., Bangkok, Thailand \,\textbullet\, $^{2}$Independent Researcher, Thailand}\par
\end{center}
\vspace{2mm}
\begin{center}
\begin{minipage}{0.92\textwidth}
\footnotesize
\noindent This document is supplementary material to the journal manuscript above; it is not a stand-alone paper. The accepted conference abstract (Appendix S0 below) remains the claim ceiling for every statement in this supplement: nothing here claims that \tphe{} improves, prevents, reduces, increases, or protects any outcome, and the extended evidence below is adjacent, analogous, or precedent evidence for constructs \tphe{} uses, never evidence that \tphe{} itself, or any of its rules, works.
\end{minipage}
\end{center}
\vspace{4mm}

\section*{Appendix S0: Accepted conference abstract (document of record)}
\footnotesize
\noindent This is the abstract accepted for the TIMA/FIMA Scientific Convention 2026 (Thailand Islamic Medical Association / Federation of Islamic Medical Associations), reproduced verbatim. It is the document of record and the claim ceiling for this manuscript: no claim in the main manuscript or elsewhere in this supplement exceeds what is stated here.

\begin{quotation}
\noindent\textbf{Background.} A family should be a space of physical, psychological, and economic safety. Whether it becomes one is shaped at the transition into marriage, a governance threshold where consent, authority, help-seeking, and access to protection are organized before a household is formed. When that threshold fails, the health consequences include violence-related injury, psychological distress, reproductive harm, and avoidable crisis-service use. Prevention frameworks such as WHO's RESPECT cover relationship skills, empowerment, safe environments, and services, yet this transition remains an underdeveloped intervention point.

\textbf{Objective.} To propose Trustworthy Premarital Health Education (T-PHE) as primary prevention infrastructure that increases each participant's control over marriage-related health decisions within safeguarding and referral boundaries.

\textbf{Methods.} This practitioner-led, document-informed conceptual synthesis drew on the governance and safeguarding procedures of ARAYA Nikah, a Muslim-led social enterprise registered in Thailand, and on its premarital education programme for Muslim and intercultural couples (over 600 participants since 2024). Recurring concerns in these sources were synthesized, through a culturally rooted ethic of limited knowledge, entrusted responsibility, human dignity, just boundaries, and repair, into a programme theory with six safeguards, a failure criterion, and outcome domains for prospective testing. Sources came from one organization; no effectiveness testing was undertaken.

\textbf{Results.} T-PHE is trustworthy only when it protects informed choice rather than relationship continuation. It rests on one mechanism: increasing each participant's control over whether to proceed with, delay, or refuse marriage and over when and how to seek help. Six safeguards protect that choice: epistemic humility, cultural intelligibility, affected-person visibility, dignity and consent, bounded authority and referral, and repairability. Educators must make power asymmetries and coercion visible, keep culture and religion from becoming unchecked authority, and refer beyond their own competence. Communication and conflict training is not universally benign. Where violence, fear, or coercive control is suspected or disclosed, joint work must stop pending confidential assessment and must not resume where coercive control is confirmed; survivor-centred support and qualified referral take precedence. T-PHE fails if it raises knowledge scores, attendance, satisfaction, or marriage completion while leaving informed choice, recognition of coercion, reproductive and mental-health literacy, safe help-seeking, and accountable referral unchanged.

\textbf{Conclusion.} Premarital health education should be recognized, governed, and tested as primary prevention infrastructure, not optional relationship advice. Public health should enter before the family becomes a site of clinical rescue.

\textbf{Keywords:} premarital health education, health promotion, intimate-partner harm, primary prevention, safeguarding
\end{quotation}

\noindent\textit{Note: accepted abstract, TIMA/FIMA 2026, document of record; the claim ceiling for the manuscript.}

\clearpage
\section{Evidence map and extended evidence review}
\label{sec:S1}
\footnotesize
This section reproduces, in full and with every citation intact, the extended literature narrative and the evidence map for propositions P1--P7 that the main manuscript's Discussion (``Comparison with the wider evidence base'') synthesises in condensed form. No evidence, citation, or finding below has been altered from the version this section was originally copied from. Unless flagged otherwise, every citation offers adjacent, analogous, or precedent evidence for a construct \tphe{} uses --- never evidence that \tphe{} itself, or any of its five rules, works.

Unless flagged otherwise, every citation below offers adjacent, analogous, or precedent evidence for a construct \tphe{} uses --- never evidence that \tphe{} itself, or any of its five rules, works; this single sentence replaces a per-subsection restatement of that caveat.

\subsection{Burden}
The population-level case for treating the marriage transition as a safeguard point rests on burden-of-disease and surveillance data, not on any evaluation of premarital education. The main manuscript's Introduction established the burden case (GBD 2023 \cite{ref1}; WHO/LSHTM \cite{ref2}: 27\% lifetime prevalence, UI 23--31\%, ever-partnered women 15--49, substantial exposure already in the 15--19 and 19--24 bands). A companion burden-of-proof analysis found a moderate association with major depressive disorder and maternal abortion/miscarriage (its largest), and only weak associations with HIV, anxiety and self-harm \cite{ref21}. A Lancet Psychiatry Commission situates these estimates within a wider call for mental-health-service investment \cite{ref22}; a Lancet Public Health review identified only 30 unique interventions from 16 countries addressing IPV and/or violence against children \cite{ref23} --- sparse relative to the burden. These are surveillance/modelling estimates, not evidence any intervention changes them.

\subsection{Primary prevention frameworks, what works, and recognition}
The strongest experimental evidence that structured relationship curricula can reduce IPV comes from community-mobilization and couples-curriculum trials outside the premarital-education literature. SASA! was associated with lower social acceptance of IPV (aRR 0.54) and a physical-IPV effect whose CI crossed 1 (aRR 0.48) \cite{ref24}, with norms change --- not knowledge or attendance --- explaining 70--95\% of the mediated effect \cite{ref25}. Indashyikirwa in Rwanda found a large reduction in women's reports of IPV (aRR 0.44) and men's reports of perpetration (aRR 0.54) \cite{ref26}, while a faith-based trial in Nigeria reported a 61\% reduction in odds of any reported IPV \cite{ref27} --- the clearest faith-setting positive result in this review. Set against these are equally rigorous null results: a scaled adaptation in rural Rwanda found no effect on any primary outcome \cite{ref28}, and a faith-leader caregiving trial in Kenya found no effect on its declared primary outcome while showing secondary IPV effects \cite{ref29} --- the outcome-substitution pattern \tphe{}'s failure criterion is designed to flag. Only 6/178 (3\%) reviews updating the WHO RESPECT framework targeted relationship-skills strengthening specifically \cite{ref30}, and a scoping review of implementation science found fidelity, acceptability and feasibility the most-studied outcomes, with scale, sustainability and survivor-centred practice minimally addressed \cite{ref31} --- a field-level gap consistent with treating \tphe{} as a programme theory rather than advancing an effectiveness claim at this stage.

Psychological inoculation --- drawn from misinformation research, not interpersonal coercion, and used here only as an analogy --- is the closest experimental literature to Layer 0 (a proposed extension)'s ``awareness as protection'' claim. The largest field experiment (7 preregistered studies, 6 RCTs, $n=6{,}464$, plus a YouTube field study $n=22{,}632$) found inoculation videos improved recognition of manipulation techniques and sharing-decision quality \cite{ref182} (effect sizes not reported in the abstract), extended to vaccine misinformation \cite{ref183}, election misinformation across 12 EU nations \cite{ref184}, a gamified format \cite{ref185}, real-world cross-protection against untrained techniques \cite{ref186}, and methodological work isolating technique from testing effects \cite{ref187} --- none of it about intimate-partner coercion. The strongest cautionary evidence comes from a 1-year cluster RCT of Coaching Boys Into Men (16 high schools, $n=1{,}513$ male athletes): perpetration fell, but recognition of abusive behaviours did not change (95\% CI $-$0.15 to 0.09) and bystander intentions were unchanged \cite{ref190} --- recognition and behaviour dissociating exactly as the failure criterion is built to catch. A Cochrane protocol for dating-violence prevention has not yet reported outcomes \cite{ref188}; a pilot questions whether attitudes proxy behaviour \cite{ref189}; further design and implementation papers report no recognition outcomes yet \cite{ref191,ref192}; a non-randomised bystander study found effects on risk recognition and perceived self-efficacy \cite{ref193}. Qualitatively, women describe difficulty naming gaslighting despite describing its effects \cite{ref194}; being silenced through negative responses to help-seeking, and loneliness, form part of the pathway to later intimate-partner violence and abuse in young adults \cite{ref195}; disclosure barriers are documented among sexual and gender-minority survivors \cite{ref196}; two studies of South and Southeast Asian women describe culturally specific understandings and barriers relevant to naming abuse \cite{ref200,ref201} (measurement tools not yet applied here: an internalised-stigma scale \cite{ref197}, Thai/general health-literacy instruments \cite{ref198,ref199}). This literature makes Layer 0's premise plausible and testable, not established; recognition and help-seeking are shown here to be separable outcomes.

\subsection{Preconception medicine and premarital screening as a proposed analogue}
Preconception medicine treats the period before conception/marriage as a distinct clinical window: the FIGO Preconception Checklist frames structured health-assessment elements for it \cite{ref11}; couple-based preconception-care programmes have been reviewed across 14 studies in 10 countries, mostly high-income, heterogeneous designs \cite{ref32}. Premarital genetic screening is the closest direct precedent for a premarital contact point at population scale (roughly 30--40\% of the Thai population carries a thalassaemia gene \cite{ref117}, motivating national couple-screening programmes): premarital sickle-cell-trait screening uptake across Africa pools under half of eligible couples (47.82\%, 95\% CI 46.53--49.11), ``suboptimal'' \cite{ref12}; premarital thalassaemia education produces consistent knowledge gains ($p<0.001$) but inconsistent screening-uptake behaviour \cite{ref13} --- the same knowledge/behaviour gap recurring below.

\subsection{Premarital/relationship education outcomes and the knowledge-to-well-being gap}
Across three cumulative meta-analyses spanning 1997--2022, the MRE evidence base has consistently measured relationship quality and communication skill \cite{ref6,ref33}, and more recently mental health and coparenting \cite{ref7}, with effects real but small ($d\approx 0.03$--$0.45$ \cite{ref6,ref7}) and authors flagging thin racial/ethnic diversity \cite{ref6} and rarely measured policy-relevant outcomes \cite{ref7}; a third meta-analysis in this lineage examined programmatic moderators such as dosage and communication emphasis \cite{ref33}. Individual RCTs complicate rather than resolve this: an 8-year follow-up found no divorce benefit and a harmful-moderation pattern for higher-conflict couples \cite{ref9}; a 3-year RCT found couples with acute baseline risk benefited less \cite{ref10}; a later meta-analysis of relationship aggression (31 studies, $n=25{,}527$) found $d\approx 0.11$ overall, not statistically convincing when restricted to RCTs ($d\approx 0.05$) \cite{ref8}. Newer, smaller trials sharpen the specific knowledge-to-well-being gap this paper labels its ``middle of the bridge'': a quasi-experimental study among engaged couples in Babol, Iran, found booklet- versus video-based premarital education improved reproductive-health-literacy access, evaluation and decision-making dimensions, but not reading-and-understanding \cite{ref34}; a 2026 study of reproductive/sexual-health education in premarital counselling classes in Rasht, Iran, found knowledge and attitudes improved at 3 months but reported no clinical outcomes \cite{ref35}; and a meta-analysis of online relationship-education programmes (study populations not characterised in the abstract) found they improved satisfaction, communication, confidence and individual-level anxiety/depressive symptoms \cite{ref36}. At the far end of this pathway, a large meta-analysis (126 articles, $>72{,}000$ participants) found marital quality associated with physical health, lower mortality risk and lower cardiovascular reactivity during conflict \cite{ref37} --- a well-being pathway exists in principle, not evidence premarital education produces the marital quality that pathway depends on. None of this literature has measured informed choice, coercion recognition, or safe help-seeking. Two frequently cited works in this literature --- Fawcett et al. (2010, \emph{Family Relations}) and Carroll \& Doherty (2003) --- are not indexed in PubMed; Fawcett et al. (2010) is cited once, flagged as not PubMed-indexed, in the novelty comparison table of the main manuscript, and Carroll \& Doherty (2003) is not cited under this review's PubMed-only scope.

\subsection{Screening, first-line response, referral, and technology-facilitated routes}
The clinical screening literature offers the clearest precedent for the distinction \tphe{}'s failure criterion is built on. A Cochrane review of 13 RCTs found screening increased identification (OR 2.95) but no evidence of an effect on referral (OR 2.24, 2 studies only) or on re-exposure at 3--18 months \cite{ref38}. The 2025 USPSTF evidence report reviewed 35 studies ($n=18{,}358$): 3 RCTs ($n=3{,}759$) comparing screening with no screening found no significant reduction in IPV or other outcomes over 3--18 months, instrument sensitivity ranged 26--87\% (specificity 80--97\%), and of 13 intervention RCTs only perinatal home-visiting and multi-risk behavioural counselling showed benefit \cite{ref39}; the companion 2025 recommendation statement retained a moderate-net-benefit judgement for reproductive-age women \cite{ref40}. Provider training may improve knowledge, attitudes and self-perceived readiness, with weak and inconsistent evidence for actual response behaviour or survivor outcomes \cite{ref41}. Where first-line response protocols have been trialled directly (WHO LIVES/ARCHES in Nigeria), results are more encouraging but bounded: reduced relative risk of IPV exposure and reproductive coercion among women not already experiencing violence, but no effect on contraceptive use \cite{ref42}, with WHO's own guidance conditioning any inquiry about violence on a functioning referral system already in place \cite{ref43,ref44}. Premarital education itself may function as a gateway rather than a substitute for later help-seeking --- one US study found it associated with an increased likelihood of later couples counselling, with `not knowing where to go' reported alongside cost as a principal barrier among low-income couples who did not seek help \cite{ref106,ref107}. Delivery agents themselves may complicate a referral-first safeguard: in a survey of Christian clergy and congregation members, clergy reported being comfortable making referrals for professional counselling, whereas most congregation members preferred counselling with their own pastor and respondents wanted more training on domestic violence --- a live tension, not settled support, for a policy that assumes faith-based delivery agents will refer rather than counsel \cite{ref118}.

Evidence on Layer 2 (a proposed extension)'s moderated peer community is genuinely two-sided. On support: a review of youth mental-health peer forums found only 2 of 3 RCTs significant, the other four null \cite{ref139}; realist syntheses across mental-health and cancer settings agree outcomes depend on moderation, trust and anonymity/privacy control, noting ``not all online support groups are helpful, and in some cases they may increase distress'' \cite{ref140,ref141,ref146}. A scoping review of moderator practice found only 5 of 397 screened articles eligible, only one using a qualified health professional as moderator \cite{ref142}; qualitative work describes the role as under-supported and under-specified \cite{ref143}, while other qualitative work documents moderators' facilitative strategies and the need for better training and supervision \cite{ref144}. Weakly moderated platforms specifically amplified harmful content relative to strongly moderated ones \cite{ref145}. A Chinese IPV-specific peer-support study reports mostly supportive, rarely blaming, content \cite{ref147}, but existing online-IPV tools focus on leaving a relationship with little post-departure attention \cite{ref148}, and digital-IPV-tool safety research remains ``virtually non-existent'' \cite{ref137}. On harm: technology-facilitated abuse maps onto every domain of the Power and Control Wheel \cite{ref149}, and more broadly spans digital dating abuse, cyberstalking, technology-facilitated sexual assault, image-based sexual abuse, and online sexual harassment \cite{ref150}; immigrant/refugee women face additional, unique barriers to escaping and reporting it \cite{ref151}; it includes systematically characterised cyberstalking \cite{ref152}; and record-keeping features in coparenting apps have been shown to be sites of both coercive control and survivor resistance \cite{ref153}. Anonymity increases self-disclosure, while other features of online communication such as asynchronicity, multiple audiences, and audience feedback favor impression management over genuine self-presentation \cite{ref154}; Muslim youth's social media use carries both cultural/religious connection and discrimination \cite{ref157}; young people's reported benefits of online help-seeking come with literacy and trust barriers \cite{ref156}. This symmetry is why Layer 2 needs a named detection mechanism, not only a withdrawal rule.

\subsection{Couple-based work, coercive control, and assessor independence}
The evidence base most directly relevant to \tphe{}'s stop rule (R4) is small and explicitly cautious. The only systematic review of couples therapy as an IPV treatment, after screening 1{,}733 citations, found only 6 studies met inclusion; pooled data suggested conjoint treatment ``viable in select situations,'' with professionals cautious about violent-retaliation risk \cite{ref45}. Clinical literature converges on screening, safety modification, ongoing reassessment, and distinguishing situational couple violence from coercive controlling violence as prerequisites for any conjoint work \cite{ref46,ref47}. Survivor-voiced evidence complicates any default toward joint response --- only 9 of 17 women in one qualitative sample rated couples counselling a good idea \cite{ref48} --- and measurement is contested: coercive control is often reported as reciprocal \cite{ref49}, standard instruments do not measure perpetration and victimisation symmetrically \cite{ref50}, and even validated disclosure instruments substantially undercount coercive control through item-interpretation confusion and fear \cite{ref51}.

No study tests \tphe{}'s no-stake assessor (a proposed extension); the closest matched evidence is regulated independence in other high-stakes consent settings, offered here as analogy, not domain evidence. Transplant medicine mandates an Independent Living Donor Advocate separated from the recipient's care team \cite{ref158,ref159}, extended by proposal to two temporally separated consent processes for domino donors \cite{ref160}; facial plastic surgery literature names the same conflict-of-interest problem without a mandated remedy \cite{ref161}, and psychotherapy literature names the same dual-relationship problem, noting only meager research and little guidance on nonsexual dual relationships \cite{ref162}, with common occurrence in small or rural settings noted elsewhere \cite{ref163}. Structured IPV risk-assessment tools show independence helps but does not guarantee accuracy: a meta-analytic review found ODARA best-performing among five compared (weighted AUC=.666), yet only 9 of 20 (45\%) predictive-validity estimates used the tool exactly as intended, with a separate study validating an authorised German translation of the ODARA instrument \cite{ref164,ref168}; the Danger Assessment and its 5-item version retain real predictive validity for severe/lethal re-assault (AUC .68--.78) \cite{ref166}; the Taiwan-adapted version shows similarly strong discriminant and predictive validity, particularly for high dangerousness, though no specific AUC value is reported in the abstract \cite{ref165}; and a version updated to account for multiple/near-loss-of-consciousness strangulation predicts near-fatal re-assault significantly better than the original Danger Assessment, though no AUC value is reported in the abstract \cite{ref167}; but multiagency structured-judgment review found it often unclear which recorded factors drove a risk categorisation \cite{ref169}; a police-administered SARA field study found intervention intensity could reduce recidivism in high-risk cases while increasing it in low-risk cases \cite{ref170}; structured-judgment ratings, not only actuarial scores, differed by ethnicity between Indigenous and non-Indigenous participants in one study \cite{ref171}. Check-in and brief-contact models show the same mixed pattern the failure criterion is designed to catch: quality-improvement warm-handoff work doubled successful staff-to-community-health-worker referrals \cite{ref172}, yet an emergency-department screening programme found only 24.5\% of screened patients recalled receiving a resource guide despite 58.5\% being screened \cite{ref173}. An SMS brief-contact trial ($n=804$) found a 22\% relative reduction in self-harm repetition at 12 and 24 months but no difference in time to first repeat event \cite{ref174}, and a caring-contacts trial found no significant difference in loneliness at 6 months by treatment arm, concluding an added phone call increased complexity without improving effectiveness \cite{ref176} (protocol for the related comparative-effectiveness trial: \cite{ref175}). Implementation-fidelity work in an adjacent prevention programme shows that measured delivery fidelity and dose do not by themselves guarantee outcome change \cite{ref179}. This literature makes assessor independence necessary but not sufficient, and supports treating the day itself as presence of the space rather than as the outcome --- consistent with the failure criterion, not a validation of it.

\subsection{Decision quality, autonomy measurement, and digital self-assessment}
Decision quality --- shared decision-making and informed, values-congruent choice --- is this paper's proposed nearest measurement family for \tphe{}'s decision-control construct (not a validated instrument for it). The ``three-talk'' model structures shared decision-making into team, option and decision talk \cite{ref15}; a Cochrane review of 209 decision-aid studies ($>107{,}000$ participants) found consistent improvement in knowledge, risk-perception accuracy and values-congruent choice, and reduced decisional conflict \cite{ref14}. The MRC complex-intervention framework \cite{ref52}, its ADAPT context-adaptation guidance \cite{ref53}, and its process-evaluation guidance \cite{ref54} supply the methodological scaffolding for describing \tphe{} as a programme theory awaiting adaptation and process evaluation, not efficacy testing. The reproductive-autonomy and coercion instruments reviewed above \cite{ref16,ref17,ref18} remain the closest content-matched but Western-derived, Thai-unvalidated, family.

Layer 0 (a proposed extension)'s anonymous scored self-check has real precedent, though the precedent does not show benefit. The most rigorously tested example, I-DECIDE, is a masked two-arm RCT ($n=422$, 79--80\% 12-month retention) built from an explicit theoretical model and published protocol \cite{ref119,ref121}, with the methodological and ethical challenges of running the web-based trial documented separately \cite{ref122}; at 6 and 12 months it found no improvement in self-efficacy, with scores in fact slightly lower in the intervention arm than control (mean difference 1.3, 95\% CI 0.3--2.3 at 6 months; 1.6, 95\% CI 0.5--2.7 at 12 months), and no between-group difference in depressive symptoms \cite{ref120}. The myPlan app family, tested across protocol, longitudinal and sub-population studies \cite{ref123,ref126,ref127,ref128}, reported a significant reduction in reproductive coercion at 12 months ($p=.019$), but its own reported $p$-value for suicide-risk improvement ($p=.46$) does not support that outcome as significant despite wording in the source abstract implying otherwise \cite{ref124}; a bystander sub-study and teen adaptation extend the programme without changing this picture \cite{ref125,ref126}. A double-blind Canadian trial of a tailored online safety and health tool (planned $n=450$ in the protocol \cite{ref129}, $n=462$ enrolled in the completed trial \cite{ref130}) is the closest design analogue to a full self-check-plus-routing model, and a user-centred design report catalogues the wider suite of screening/safety apps this field has produced \cite{ref131}. The only Thailand-based example, a Thai adaptation of iCanPlan, has so far been evaluated by experts only (usability 4.93/5, cultural appropriateness 5.00/5), not by users \cite{ref132}; a military-partner survey and two EHR/computer-assisted screening trials extend the pattern without adding a premarital, scored, self-directed instrument \cite{ref133,ref134,ref135}. A Kenya-adapted risk instrument shows a scored tool can be validly tailored to a non-Western setting (AUC 0.78) \cite{ref136}, but a systematic review of ICT-based IPV interventions notes safety, equity and unintended-consequence research is ``virtually non-existent'' \cite{ref137}, and a review of web-based/mHealth IPV interventions finds computer-based/self-administered screening is the most common, not the best-proven, approach \cite{ref138}. This literature positions Layer 0 as feasible and acceptable, not validated: no record here is evidence a premarital scored self-check changes recognition, help-route knowledge or enacted choice.

\subsection{Multicultural clinical safety, language access, and Islamic medical ethics}
This is the evidence most directly relevant to \tphe{}'s cultural-intelligibility safeguard, and it is consistently cautionary. Cultural-competence training for health professionals repeatedly shows provider-level knowledge or attitude gains without a corresponding patient-outcome gain: a Cochrane review of 5 RCTs (337 professionals, 8{,}400 patients) rated the evidence low-quality, with some improvement in involvement/understanding but no reliable clinical-outcome effect, and adverse outcomes not assessed in any trial \cite{ref55}; a narrower review (7{,}879 records to 5 outcome-reporting studies) found no significant patient-outcome improvement in any \cite{ref56}; and Indigenous cultural-safety training across four high-income countries found outcomes overwhelmingly provider-level, with weak community-involvement and clinical-outcome evidence \cite{ref57}. Cultural humility \cite{ref58} and cultural safety --- a shift ``from provider control to patient agency'' \cite{ref59} --- are the constructs this paper adopts for the cultural-intelligibility safeguard, given the pattern above in which content-delivery training alone has not been shown to close this outcome gap. Language-access evidence sharpens the case for interpretation as distinct from cultural content: language discordance is associated with worse outcomes in over half of studies reviewed (84{,}750 participants) \cite{ref60} and with higher unplanned readmission/ED revisit \cite{ref61}; professional interpretation improves patient satisfaction and communication quality relative to ad hoc family/friend interpretation \cite{ref62}, and produces fewer linguistic errors than ad hoc interpretation \cite{ref63}; and a Canadian cohort ($n=124{,}583$) found language-concordant care associated with 36\% lower major cardiovascular events among allophone speakers \cite{ref64}. Thailand-specific evidence situates this in \tphe{}'s delivery context: interpretation/cultural-mediation for migrant health in Thailand faces budget, training and system-status constraints \cite{ref65}; migrant-health coverage in Thailand is under-studied in reproductive/maternal domains relative to general health \cite{ref66}; a Thai maternity-care ethnographic study named ``the system is in control,'' with insensitivity to cultural practice a subtheme \cite{ref67}; Thai Muslim NCD patients in the rural south report medication-adherence barriers (including Ramadan-related non-adherence) tied to social/cultural context \cite{ref68}; and Muslim mothers describe faith, patience and family/professional support as their coping structure after their children's cardiac surgery \cite{ref69}. Migrant/Muslim-women SRH-access reviews consistently name information, language and cultural-difference barriers \cite{ref70,ref20}, while quasi-experimental tailored delivery in southern Thailand nearly doubled cervical-screening uptake \cite{ref71} --- a single behavioural endpoint, not evidence that tailoring transfers to \tphe{}'s decision-control outcome.

Tier-one Islamic bioethics literature gives \tphe{}'s ethical frame real companions in the medical literature, without becoming evidence the frame, or \tphe{}, works. Islamic ethico-legal reasoning has been mapped onto the four principles of biomedical ethics \cite{ref225} and surveyed as a field spanning virtue ethics to jurisprudential rulings \cite{ref228,ref230}; empirical work with Muslim clinicians describes real moral distress reconciling Islamic and secular institutional ethics and calls for religious-competence training and peer support \cite{ref227}, while cross-gender-interaction norms are addressed separately \cite{ref226}, and end-of-life analyses treat consent as a shared, family-involved process that does not displace individual assent \cite{ref229}. Two Journal of IMA records, admitted under a labelled professional-association carve-out rather than as leading-journal evidence, frame informed consent as properly incorporating religious/cultural values \cite{ref241} and frame Qur'an-guided physician character as trustworthy, accountable professional conduct \cite{ref242}. Faith-integrated delivery can carry a measurable clinical effect: a cluster-randomised mosque-based diabetes-prevention trial reduced 12-month type 2 diabetes incidence from 17.1\% to 9.8\% (adjusted hazard ratio 0.75, 95\% CI 0.60--0.95, $P=.02$) \cite{ref231}, and a mosque-based, lay-led Islamic Trauma Healing programme for refugee PTSD showed effect sizes of $d=-0.67$ (PTSD severity), $d=-0.66$ (depression) and $d=0.71$ (well-being), maintained at 12-week follow-up \cite{ref232} --- though two published errata (\emph{JAMA Netw Open} 2024;7(9):e2440198 and 2024;7(10):e2446257) mean the corrected article, not the original abstract alone, should be checked before any number from it is repeated. A religiosity-anxiety meta-analysis found a consistent protective association (mean $r=-0.22$ across 10 studies) \cite{ref238} (further tier-one companions on vaccination decision-making, critical-illness family needs and faith-based infectious-disease framing: \cite{ref233,ref234,ref243}). This literature also complicates any simply-protective reading: a systematic review of gender-based violence among migrant Muslim women found the husband's control over the woman's life a risk factor for violence at the microsystem level, even where spirituality was protective at the individual level \cite{ref235}, and a systematic review of self-immolation found marital and familial conflict a principal driver in some Muslim-majority settings \cite{ref239} (further tier-one companions on sexual/reproductive health, assisted reproduction and religion-and-health: \cite{ref236,ref237,ref240}). No indexed premarital guideline from the Federation of Islamic Medical Associations (FIMA) or the Islamic Medical Association of North America (IMANA), and no Thai Muslim premarital trial, exist in this record. This is what the abstract's ethical terms are declared alongside, in their own field's own words, with a plain gloss attached --- the frame's companion literature, never evidence that the frame prevents harm.

\subsection{Thailand and the southern border provinces}
Direct Thai evidence remains sparse and geographically uneven. A four-region national survey ($n=2{,}462$ married/cohabiting women) found roughly 15\% lifetime psychological/physical/sexual IPV and controlling behaviours ranging 4.6--29.3\% \cite{ref72}; clinic-based studies found 21\% past-year IPV across two Bangkok gynaecology clinics (IPV associated with STI, aOR 2.65) \cite{ref73} and 13.1\% ever-abused among 475 pregnant Thai women \cite{ref74}. Provider-side evidence --- from \emph{northeast} Thailand, not the southern border --- found providers valued IPV screening but cited cultural-norm, training and service-condition barriers \cite{ref75}, a caution against generalizing provider attitudes across Thai regions. Southern-border-specific evidence is more extensive: among 700 Muslim adolescents on the southern border, 9\% were sexually experienced, with condom non-use, STI and repeat pregnancy concentrated in that subgroup \cite{ref76}; among 360 Muslim army conscripts in the three deep southern provinces, 78.7\% showed poor HIV/STI-transmission knowledge, examined further through an Islamic-values framing of sexual knowledge in a related conscript sample \cite{ref77,ref78}. Mental-health evidence from the conflict-affected southernmost provinces is consistent in direction across cohorts: migration and economic stress were associated with more psychiatric symptoms and religiosity with fewer among 2{,}053 Muslim adults \cite{ref79}; among youth 18--24, parental absence and travel-safety worry were associated with more psychiatric symptoms, with religious practice protective for men \cite{ref80}; among 387 adolescents, 7.2\% screened at risk of depression, with family dysfunction and PTSD strongly associated with depression risk \cite{ref81}; and young adults wishing to migrate for work reported more psychiatric symptoms and lower happiness, with chronic violence linked to education, employment and settlement plans \cite{ref82}. The same religious-interpretation and trust structure recurs in vaccine-acceptance evidence: per interviewed health workers, perceived religious prohibition, spousal disagreement, fear of side effects and religious-leader views were the leading barriers to MMR-vaccine acceptance among Muslim parents \cite{ref83}, and childhood-immunization beliefs were shaped by self-interpretation, misinformation and halal concerns, tempered by reliable information and community health-network leaders \cite{ref84}. No study identified in either round evaluates a premarital-education programme, of any design, with a Thai Muslim or Thai intercultural population.

\subsection{Culture as a health determinant, humility, and self-regulation}
Framed at the widest level, culture-and-health scholarship supplies theoretical warrant, not efficacy evidence, for treating culture as a determinant \tphe{}'s safeguards must hold rather than resolve. The Lancet Commission on Culture and Health argues culture shapes the meaning of health, behaviour, trust and communities of care \cite{ref85}; a companion Lancet series maps mechanisms --- separation and hierarchical power \cite{ref86}, and pathways through health systems and communities \cite{ref87} --- by which racism and discrimination affect health, and a 2026 follow-up renews calls for inclusive services, equitable access, rights and dignity for people on the move \cite{ref88}. Evidence on whether \emph{adapting} an intervention to a specific culture adds benefit beyond a generically delivered one is more equivocal than commonly assumed: a meta-analysis of 67 RCTs of culturally adapted psychological interventions for people of Chinese descent found a medium overall effect but \emph{no significant difference} between culturally modified and culturally specific adaptation types \cite{ref89} --- evidence against assuming any particular adaptation strategy is superior, not against cultural adaptation itself. Two structured-social-support RCTs outside the marriage-transition setting offer an analogy, not evidence, for the mechanism \tphe{}'s safeguards are built around: a Nepal social-capital-building programme was associated with lower incident major depressive disorder (7.4\% vs 12.9\%) \cite{ref90}, and a UK family-wellbeing RCT (674 participants, 62\% ethnic-minority) improved parental mental wellbeing post-intervention and at 6 months \cite{ref91}. The pattern is read here as suggestive that structured social support, not content delivery alone, may be the active ingredient --- a hypothesis adopted by analogy, not one either trial tested directly.

The constructs the model uses to gloss \emph{faqr} and \emph{taqwa} --- humility, compassion and self-regulation --- are real, measured psychological variables, but the evidence sits mostly in general and clinical populations, not premarital or Muslim ones, and the sequence to safety is not self-executing. Intellectual humility predicts more constructive responses to conflict \cite{ref203}, greater empathic accuracy toward outgroup members \cite{ref204}, and greater adherence to expert recommendations even controlling for political affiliation \cite{ref205}; cultural humility has a developed theoretical basis \cite{ref206,ref207} and institutional training improves attitudes and knowledge, with the authors themselves cautioning that a single training session is insufficient \cite{ref208}. Self-compassion has a large protective association with psychopathology ($r=-0.54$, 95\% CI $-0.57$ to $-0.51$) \cite{ref209}, and compassion-focused therapy improves self-compassion and reduces symptoms in trial evidence \cite{ref210}, but its relational benefit is conditional: among men, self-compassion predicted constructive repair only when conscientiousness was high, and the association reversed at low conscientiousness \cite{ref211}; two further studies link self-compassion to relationship satisfaction and conflict resolution in general and Black heterosexual-couple samples \cite{ref212,ref213}. Empathy shows the same double edge: working memory was associated with less physical abuse perpetration, an effect mediated by affective empathy \cite{ref214}, but in one clinical sample of men, emotional empathy was positively associated with perpetrated psychological IPV \cite{ref215} --- restraint in the use of power, not empathy alone, is the load-bearing step. Self-regulation deficits (inattention, impulsivity) are cohort-study risk factors for adult IPV perpetration, and existing treatment programmes are reported to have minimal impact to date \cite{ref216}; a systematic review traces the broader psychosocial pathway from childhood intra-familial victimisation to adult relationship violence \cite{ref217}. Religiosity and spirituality were protective for general interpersonal aggression but \emph{not} specifically for domestic violence in a 67-study systematic review and meta-analysis \cite{ref220}, complicating any simple ``faith protects'' reading (further tier-one companions extending without resolving this picture: \cite{ref221,ref223}). Trust and dyadic-maintenance scales \cite{ref218,ref219}, a humility-forgiveness mediation study \cite{ref202}, and the only premarital analogues in this literature --- an 8-couple feasibility pilot of a culturally adapted couple-learning programme \cite{ref100} and a premarital HIV-counselling cluster RCT \cite{ref224} --- round out the picture. Humility and compassion are evidence-bearing constructs by analogy from therapy and general populations; they position R1's ethical language, they do not supply evidence that R1 works in this setting.

\begin{table}[H]
\centering\footnotesize
\caption{Evidence map for propositions P1--P4 (reference numbers as in this supplement's own bibliography, which is numbered independently of the main manuscript; each entry carries its PMID). Continued in Table~\ref{tab:evidencemap2} for P5--P7.}
\label{tab:evidencemap}
\begin{tabularx}{\textwidth}{@{}p{0.9cm}YYY@{}}
\toprule
& \textbf{Relevant / complicating records} & \textbf{Evidence level} & \textbf{What remains untested}\\
\midrule
P1 & \emph{Relevant:} burden/instrument base \cite{ref1,ref2,ref5,ref16,ref18,ref92}; norms (not knowledge) drove 70--95\% of the SASA! effect \cite{ref26,ref25}; decision-aid/shared-decision-making analogue \cite{ref14,ref15}. \emph{Complicating:} RCT-only aggression effect not convincing \cite{ref6,ref33,ref7,ref8}; uneven knowledge/literacy-to-outcome transfer \cite{ref34,ref35}; Muslim identity does not uniformly predict lower autonomy \cite{ref93}. & Mixed: burden data and single-time-point instrument-validity studies are strong; the causal claim that a governance layer \emph{adds} decision control beyond content, and that decision-quality measures capture it, is untested for any premarital-education format. & Whether a governance/safeguard layer added to premarital content produces a measurable decision-control effect distinct from content alone, using any instrument (existing or adapted), has not been tested.\\
\addlinespace
P2 & \emph{Relevant (adjacent precedent only --- no cited record tests a private channel specifically inside premarital/marriage education):} \cite{ref94,ref95,ref96,ref97,ref98,ref99}. \emph{Complicating:} even validated instruments substantially undercount coercive control \cite{ref51}. & Mixed: instrument-validation and single trials; retention/attrition data are a direct caution. & Whether a private channel embedded specifically in premarital education (not a clinical encounter) reliably surfaces concealed coercion has not been tested.\\
\addlinespace
P3 & \emph{Relevant:} \cite{ref100,ref101,ref27}; tailored edutainment nearly doubled cervical-screening uptake, a single behavioural endpoint \cite{ref71}. \emph{Complicating:} \cite{ref102,ref103,ref19}; cultural-competence training repeatedly improves provider knowledge/attitude without patient-outcome gain \cite{ref56,ref57,ref55}; no efficacy difference between adaptation types in one large meta-analysis \cite{ref89}. & Mixed to low: mostly single small trials, pilot studies, and systematic reviews of provider-level rather than patient-level outcomes; no controlled comparison of tailored vs untailored delivery on a safety outcome. & Whether cultural tailoring can be operationalized so that it improves engagement without ever being permitted to override a consent/safety threshold has not been tested in any premarital-education trial.\\
\addlinespace
P4\slash STOP & \emph{Relevant/analogous (not effectiveness support):} \cite{ref45,ref46,ref47,ref48,ref9,ref10,ref43}. \emph{Complicating:} \cite{ref49,ref50}. & Bounded: the strongest single review here is a meta-analysis of only 6 small studies. & No study has tested \tphe{}'s specific stop rule (R4) against a graduated-response alternative; the assessment mechanism that would trigger it is not specified or validated in \tphe{}'s own materials.\\
\bottomrule
\end{tabularx}
\end{table}

\begin{table}[H]
\centering\footnotesize
\caption{Evidence map for propositions P5--P7, continued from Table~\ref{tab:evidencemap}.}
\label{tab:evidencemap2}
\begin{tabularx}{\textwidth}{@{}p{0.9cm}YYY@{}}
\toprule
& \textbf{Relevant / complicating records} & \textbf{Evidence level} & \textbf{What remains untested}\\
\midrule
P5 & \emph{Relevant:} 2025 USPSTF evidence report and recommendation statement \cite{ref38,ref39,ref40}; referral pathway \cite{ref42,ref43,ref44,ref104,ref105,ref106,ref107}; language-concordant referral \cite{ref60,ref62,ref63,ref61,ref64}. \emph{Complicating:} provider training improves knowledge/readiness, not reliably behaviour \cite{ref108,ref109,ref41}. & Strong for the general principle; weak/absent for \tphe{}'s own referral pathway, and untested for whether it meets language/cultural-accessibility standards in a Thai setting. & Whether \tphe{}'s proposed referral pathway meets timeliness/accessibility/follow-through \emph{and} language/cultural-accessibility thresholds in a Thai setting has not been tested.\\
\addlinespace
P6 & \emph{Relevant:} \cite{ref110,ref28,ref111,ref112}; a faith-leader trial found no effect on its primary outcome (secondary IPV effects only) \cite{ref29}; adjacent-field confirmation: screening uptake ``suboptimal'' despite an established programme \cite{ref12}; knowledge gains consistent, uptake inconsistent \cite{ref13}; RHL gains uneven across dimensions \cite{ref34}; knowledge/attitudes rose, no clinical outcomes reported \cite{ref35}; provider training raises knowledge/attitude without patient-outcome gain \cite{ref55,ref56,ref57}. & Strong: multiple RCTs and systematic reviews across two literatures now directly demonstrate this knowledge/behaviour and process/outcome dissociation. & Whether this dissociation pattern specifically occurs in \tphe{}'s own programme has not been tested --- it is now well-grounded across two adjacent literatures, not yet observed in \tphe{}.\\
\addlinespace
P7 & \emph{Relevant:} \cite{ref113,ref114,ref31,ref115,ref116}; methodological scaffolding for future revision cycles: MRC complex-intervention framework \cite{ref52}, ADAPT guidance \cite{ref53}, MRC process-evaluation guidance \cite{ref54}. & Weak to moderate: mostly programme-description and narrative-review evidence, plus general-purpose methodological guidance, not outcome trials of repair mechanisms themselves. & Whether \tphe{} has, or will build, an actual mechanism (audit cycle, feedback loop, revision trigger) for repairability has not been specified or tested.\\
\bottomrule
\end{tabularx}
\end{table}

\subsection{Gap statement}
\tphe{} is offered as a programme theory --- a structured account of the mechanism by which a governance layer of five rules (R1--R5), including a stop--assess--refer--record rule (R4), is hypothesized to preserve or restore each participant's decision control at the marriage transition --- and not as a validated or effective intervention. No component of \tphe{} has been tested, in any population, against the outcomes its own failure criterion names as necessary (informed choice, coercion recognition, reproductive and mental-health literacy, safe help-seeking, accountable referral). This review adds two findings that sharpen, without resolving, the case for testing: first, that a premarital contact point can plausibly be built and delivered at population scale (drawing on preconception medicine and premarital genetic screening as the nearest analogue, not a tested precedent for \tphe{} itself), while knowledge gained at that contact point does not reliably convert into the behavioural or decision outcome desired, in either the screening or the relationship-education literatures; second, that cultural-competence and cultural-tailoring interventions elsewhere in health care repeatedly raise provider- or participant-level knowledge and attitudes without a demonstrated patient-outcome gain, and that no adaptation strategy has been shown superior to another in the one large trial that tested this directly. No existing instrument for any of \tphe{}'s target outcomes has been validated in a Thai or Thai Muslim population, and the decision quality/shared-decision-making measurement family this paper proposes as the nearest analogue for ``decision control'' is itself an unvalidated adaptation, not a borrowed instrument. \tphe{} should be read, and should be tested, as a falsifiable synthesis awaiting empirical evaluation --- never cited, by its authors or by others, as evidence that premarital education of this design improves, prevents, or reduces any outcome.

\section{Illustrative scenario modelling (not evidence)}
\label{sec:S2}
\footnotesize
This section reproduces the person-level Monte Carlo scenario simulation reported in the main manuscript (How we developed the framework: Analysis). It is a sensitivity exercise over declared assumptions, not empirical data about any delivered programme, and it is reported here in full because the main text carries only a one-sentence summary. \textbf{All figures in this section are illustrative scenario outputs of the model's own declared assumptions, not measured effect sizes or effect estimates for \tphe{}; they show what the assumptions coded into the model imply, offered as input to co-design.}

A person-level Monte Carlo simulation was run under declared assumptions --- two independent risk
draws per couple, classroom-level facts logged once per classroom, a single assessor draw per case,
and pre-arranged versus day-of logistics as one switch --- at 100 classrooms (eight couples each, five
seeds) and again at 1,000 classrooms (three seeds). This is a result of the assumptions coded into
the model, not empirical data about any delivered programme: no participant was observed, and nothing
here is an effectiveness estimate. It is reported because it shows which operator-controllable
conditions the model's own assumptions make consequential, checked by an independent adversarial
re-run across additional seeds.

That re-run confirmed the rankings from the original seeds while showing the first point estimates
were seed-specific and should be reported as ranges. Across seeds 1--8 and both scales, in the
simulation, under its assumptions, pre-arranging an assessor with no stake in the wedding, a private
room, and a professional interpreter \emph{before} the day was the single largest lever: true-risk
cases handled safely rose $\approx\!2.1$--$2.5\times$, cases left stuck with no assessor fell
$2.5$--$4.9\times$, and private slots overheard by family fell $6$--$10\times$. These rankings held
across seeds and at ten-fold scale; the magnitudes did not, hence the ranges.

Two further conditions followed. In the simulation, under its assumptions, once logistics were
pre-arranged, untrained facilitators became the dominant residual cause of unsafe continuation. In the
simulation, under its assumptions, a narrow, trained trigger list reduced false triggers relative to a
broad one. Day length behaved differently: in the simulation, under its assumptions, walk-outs
($\approx$5\% of couples, modelled) were unchanged by any safeguarding lever and moved only when the
dose --- the length of a single day --- was redesigned, for example by moving content to a follow-up
slot. Table~\ref{tab:delivery} summarises the four conditions the output ranked as decisive for
delivery, as distinct from the model elements themselves. None of this is a finding about \tphe{} in
practice; it states what the declared assumptions in the simulation code imply, offered as input to
co-design, not as evidence that any condition improves an outcome.

\begin{table}[H]
\centering\footnotesize\raggedright
\caption{Four delivery conditions the simulation output ranked as decisive (not model elements;
output of assumptions, not data). \textbf{All figures below are illustrative scenario outputs of the
model's own declared assumptions, not measured effect sizes or effect estimates for \tphe{}.} The
assessor row's ``no stake'' condition is a proposed extension supported only by
analogy evidence (\protect\cite{ref158,ref159,ref162,ref163}; see the main manuscript's The framework, Table~1).
A case the model counts as ``stuck'' is one where the R4 referral pathway --- to a qualified
multidisciplinary service such as Thailand's One-Stop Crisis Centre (OSCC) or an equivalent
designated service --- could not be completed in the simulation's own assumptions, not one in which
OSCC access itself was tested.}
\label{tab:delivery}

\noindent\textit{Illustrative scenario outputs, not effect sizes: figures below restate what the model's own declared assumptions imply, not a measured or estimated effect of \tphe{}.}\smallskip

\begin{tabularx}{\columnwidth}{@{}p{2.4cm}Y@{}}
\toprule
\textbf{Condition} & \textbf{What the output attributes to it, in the simulation}\\
\midrule
Assessor, room, and interpreter pre-arranged before the day & Largest lever modelled: safely handled
cases $\times2.1$--$2.5$; stuck cases $\div2.5$--$4.9$; overheard private slots $\div6$--$10$.\\
Facilitator training gate before delivery & Once logistics are pre-arranged, the residual cause of
unsafe continuation in the model.\\
Narrow, trained trigger list & Fewer false triggers than a broad or untrained list, in the model.\\
Follow-up slot rather than a longer single day & The only lever in the model that moved walk-outs,
which safeguarding levers alone did not change.\\
\bottomrule
\end{tabularx}
\end{table}

\section{Candidate outcome and measurement matrix, and claim--evidence ledger}
\label{sec:S3}
\footnotesize
\subsection{Candidate outcome and measurement matrix}
\begin{table}[H]
\centering\footnotesize
\caption{Candidate outcome domains, indicators, data sources, timing, and safety notes for prospective testing. None of these has yet been measured against an outcome.}
\label{tab:S3a}
\begin{tabularx}{\textwidth}{@{}p{3cm}YYp{2.2cm}Y@{}}
\toprule
\textbf{Domain} & \textbf{Candidate indicator} & \textbf{Data source} & \textbf{Timing} & \textbf{Safety note}\\
\midrule
Decision control & Adapted decision-quality\slash shared-decision-making instrument, context-adapted per the MRC ADAPT guidance \cite{ref53}; no instrument yet validated in a Thai or Thai Muslim population. & Individual confidential survey & Baseline, post, follow-up & Avoid partner-visible completion\\
Option awareness & Correct identification of proceed\slash pause\slash delay\slash refuse\slash help options & Knowledge task & Post & Does not prove enactability\\
Private voice & Safely offered; perceived freedom to speak & Process log, confidential survey & During, post & Do not record disclosure content unnecessarily\\
Coercion recognition & Scenario-based recognition & Individual assessment & Baseline, post & Trauma-informed fictional scenarios\\
Mental-health literacy & Recognition and help-route knowledge & Survey, scenario & Baseline, post & Not diagnosis\\
Reproductive-health literacy & Consent, contraception, reproductive coercion, care pathways & Survey, scenario & Baseline, post & Protect sexual-health privacy\\
Legal-protection literacy & Status verification, documentation, available help & Survey, task & Baseline, post & Jurisdiction-specific\\
Stop-rule fidelity & Correct facilitator response in simulation or audit & Observation & Training, delivery & Independent review preferred\\
Referral fidelity & Offer, consent, connection, latency, failed route & Safeguarding log & Trigger, follow-up & Minimum necessary data\\
Cultural intelligibility & Comprehension and relevance across groups & Survey, interview & Post & Analyse differential outcomes\\
Repairability & Appropriate review or correction completed & Audit & Ongoing & Preserve protected audit trail\\
Harms & Privacy breach, retaliation concern, distress, unsafe referral, inappropriate restart & Incident system, safe follow-up & Immediate, ongoing & Independent safety governance\\
Conventional outcomes & Attendance, satisfaction, completion, knowledge & Programme records & Post & Secondary; never sufficient alone\\
\bottomrule
\end{tabularx}
\end{table}

\subsection{Claim--evidence ledger}
\begin{table}[H]
\centering\footnotesize
\caption{Claim--evidence ledger: what evidence supports each load-bearing claim, what wording is currently allowed, and what is not yet allowed.}
\label{tab:S3b}
\begin{tabularx}{\textwidth}{@{}YYYYY@{}}
\toprule
\textbf{Claim} & \textbf{Evidence type} & \textbf{Current support} & \textbf{Allowed wording} & \textbf{Not yet allowed}\\
\midrule
IPV produces major health burden & Externally observed or modelled & WHO and GBD sources & ``IPV is a major public-health problem with physical, mental, sexual, and reproductive consequences.'' & Local causal estimates not measured\\
Thailand recorded an average 42 affected persons per day in 2024 & Official administrative report & DWF page & ``Recorded cases averaged 42 persons per day across children, women, or family violence.'' & ``Thai prevalence is 42/day''; Muslim-specific claim\\
ARAYA reached 600+ participants since 2024 & Organisation-reported practice context & Internal programme record & ``Programme reach exceeded 600 participants.'' & Effectiveness, representativeness, harm reduction\\
Decision control is the \tphe{} mechanism & Conceptual derivation & This synthesis & ``\tphe{} proposes decision control as its mechanism.'' & ``Decision control has been validated as the mechanism.''\\
Six safeguards protect decision control & Conceptual derivation & This synthesis and practice-source logic & ``Six safeguards are proposed.'' & ``The safeguards are proven effective.''\\
Stop rule is safer when fear or coercion appears & Normative design and external clinical safety logic & WHO guidance; bounded couple-intervention literature & ``\tphe{} requires routine joint work to stop pending qualified assessment.'' & Comparative safety superiority before testing\\
\tphe{} improves outcomes & Open empirical question & None & ``\tphe{} is designed to improve\ldots'' / ``should be tested for\ldots'' & ``\tphe{} improves,'' ``prevents,'' ``reduces''\\
Framework transfers beyond ARAYA & Open empirical question & Conceptual relevance only & ``Potentially adaptable; requires local testing.'' & Universal applicability\\
\bottomrule
\end{tabularx}
\smallskip

\noindent\footnotesize\textit{Note.} The six safeguards named in the row above are operationalised in the main manuscript (The framework) as five rules, R1--R5: epistemic humility and cultural intelligibility are jointly implemented as R1; affected-person visibility maps to R2; dignity and consent to R3; bounded authority and referral splits across R1 and R4; repairability maps to R5. Every safeguard survives; ``six safeguards'' and ``five rules'' name the same commitments at two levels of description throughout this work.
\end{table}

\section{Operational checklist: minimum viable \tphe{} pathway}
\label{sec:S4}
\footnotesize
\textbf{Before delivery.} Define scope and exclusions; verify referral services and emergency routes; train and assess facilitators; name a safeguarding lead, independent of routine facilitation duty, who owns supervision, escalation decisions and incident review; establish supervision and escalation; prepare language access and non-Thai pathways; separate administrative and proposed safeguarding record systems, with access restricted to the named safeguarding lead and assessor, a defined retention period, and a documented protocol for third-party requests (including from a spouse, family member or religious authority) that defaults to non-disclosure absent a legal requirement; explain consent, privacy, and limits of confidentiality.\\
\textbf{During routine education.} Mark source and uncertainty; make content culturally intelligible; offer each person a safely designed private voice; reconfirm the right to pause, delay, refuse, withdraw, or seek help; observe for fear, threats, control, unsafe participation, or out-of-scope needs.\\
\textbf{When triggered.} (0)~If danger appears imminent, the routine sequence below is suspended in favour of the organisation's emergency protocol: immediate safety planning and contact with emergency services or the OSCC emergency line take precedence over scheduling confidential assessment; facilitators are trained to recognise the difference and are never the final decision-maker for either pathway. (1)~Stop routine joint work. (2)~Protect immediate privacy and safety. (3)~Arrange confidential assessment within defined competence. (4)~Refer to qualified services. (5)~Record the minimum necessary reason, action, referral, and safe follow-up. (6)~Review; do not restart joint work without documented safe assessment. (7)~Do not resume couple mediation where coercive control is confirmed.\\
\textbf{After delivery.} Review referral connection and barriers where safe and consented; correct records or decisions when new information emerges; audit fidelity and harms; use incidents and complaints for system repair.

\subsection{Co-design governance commitments (first phase)}
A first co-design phase (main manuscript, How to test it with us) would be chaired independently of ARAYA Nikah, with decision authority in a steering group on which the organisation holds one seat and no veto. Materials co-developed there would be held by the co-design consortium under a Creative Commons Attribution (CC~BY~4.0) licence, not by the company, and partners retain the right to withdraw at any phase.

\section{Construct dictionary}
\label{sec:S5}
\footnotesize
\begin{table}[H]
\centering\footnotesize
\caption{Construct dictionary: working definitions, what each construct is not equivalent to, its observable trace, and its current status.}
\label{tab:S5}
\begin{tabularx}{\textwidth}{@{}p{2.6cm}YYYp{2.6cm}@{}}
\toprule
\textbf{Construct} & \textbf{Working definition} & \textbf{Not equivalent to} & \textbf{Observable trace} & \textbf{Status}\\
\midrule
Decision control & Usable, timely, safe capacity to proceed, pause, delay, refuse, withdraw, revise, or seek help & Satisfaction; knowledge alone; relationship continuation & Options explained, private voice, enacted pause, safe help route & Proposed primary mechanism; requires validation\\
Epistemic humility & Admission of limits, source transparency, consultation, accountable authority & Indecision; passivity; silence & Uncertainty marked, roles stated, out-of-scope consultation & Conceptually derived\\
Cultural intelligibility & Meaningful guidance within culture without unchecked cultural authority & Cultural competence as mastery; cultural exemption from safety & Language adaptation, source labelling, comprehension and enactability check & Conceptually derived\\
Affected-person visibility & Direct, safe visibility of each person who bears consequences & Couple treated as one voice; forced disclosure & Routine private route, safe contact, privacy controls & Conceptually derived\\
Dignity and consent & Ongoing respect for each person's decision and refusal & One-time signature; family or programme permission & Consent revisited; pause or refusal enacted without penalty & Conceptually derived\\
Bounded authority and referral & Role limits connected to competent, timely service pathways & Abandonment; referral leaflet alone & Stop criteria, service map, warm referral, follow-through & Conceptually derived\\
Repairability & Review and correction when new information appears & Erasure of audit trail; automatic restart & Review points, corrected record, incident learning & Conceptually derived\\
Trustworthiness & Earned property of a governed programme whose observable practice protects choice and safety & Branding; participant trust alone & Safeguard fidelity, harms, accountable referral, failure criterion & Higher-order evaluative construct\\
\bottomrule
\end{tabularx}
\end{table}

\clearpage
\section*{Supplementary references}
\noindent\footnotesize\textit{Note: reference numbers in this supplement are independent of, and do not match, the numbering in the main manuscript's reference list.}\smallskip

\begin{thebibliography}{999}\footnotesize
\bibitem{ref1} GBD 2023 IPV and SVAC Collaborators. Disease burden attributable to intimate partner violence against females and sexual violence against children in 204 countries and territories, 1990-2023: a systematic analysis for the Global Burden of Disease Study 2023. Lancet (London, England). 2026. PMID: 41386261.
\bibitem{ref2} Sardinha L. Global, regional, and national prevalence estimates of physical or sexual, or both, intimate partner violence against women in 2018. Lancet (London, England). 2022. PMID: 35182472.
\bibitem{ref5} Islam MK. Attitudes to and experiences of intimate partner violence among Rohingya women who married before eighteen years of age. Global Health Action. 2021. PMID: 34323166.
\bibitem{ref6} Hawkins AJ. Does marriage and relationship education work? A meta-analytic study. Journal of Consulting and Clinical Psychology. 2008. PMID: 18837590.
\bibitem{ref7} Hawkins AJ. How effective are ACF-funded couple relationship education programs? A meta-analytic study. Family Process. 2022. PMID: 35040124.
\bibitem{ref8} Karantzas Gery C. Do relationship education programs reduce relationship aggression? A meta-analytic study. Clinical Psychology Review. 2023. PMID: 37499317.
\bibitem{ref9} Markman HJ. A randomized clinical trial of the effectiveness of premarital intervention: moderators of divorce outcomes. Journal of Family Psychology. 2013. PMID: 23421844.
\bibitem{ref10} Williamson HC. Risk moderates the outcome of relationship education: a randomized controlled trial. Journal of Consulting and Clinical Psychology. 2015. PMID: 25664643.
\bibitem{ref11} Benedetto Chiara. FIGO Preconception Checklist: Preconception care for mother and baby. International Journal of Gynaecology and Obstetrics. 2024. PMID: 38426290.
\bibitem{ref12} Dilli Priscilla Peter. Update on the practice of premarital screening for sickle cell traits in Africa: a systematic review and meta-analysis. BMC Public Health. 2024. PMID: 38822327.
\bibitem{ref13} Anwar. Characteristics and Effectiveness of Premarital Thalassaemia Educational Interventions Worldwide: A scoping review. Sultan Qaboos University Medical Journal. 2025. PMCID: PMC12445309; not PubMed-indexed.
\bibitem{ref14} Stacey Dawn. Decision aids for people facing health treatment or screening decisions. Cochrane Database of Systematic Reviews. 2024. PMID: 38284415.
\bibitem{ref15} Elwyn G. A three-talk model for shared decision making: multistage consultation process. BMJ. 2017. PMID: 29109079.
\bibitem{ref16} Upadhyay Ushma D. Development and validation of a reproductive autonomy scale. Studies in Family Planning. 2014. PMID: 24615573.
\bibitem{ref17} McCauley Heather L. Psychometric properties and refinement of the Reproductive Coercion Scale. Contraception. 2017. PMID: 27639927.
\bibitem{ref18} Upadhyay Ushma D. Development and Validation of the Sexual and Reproductive Empowerment Scale for Adolescents and Young Adults. The Journal of Adolescent Health. 2021. PMID: 32690468.
\bibitem{ref19} Dagher Myriam. Adaptation and psychometric assessment of a sexual and reproductive empowerment scale in Arabic among refugee and non-refugee adolescent girls. BMC Medical Research Methodology. 2024. PMID: 39266993.
\bibitem{ref20} Alomair N. Factors influencing sexual and reproductive health of Muslim women: a systematic review. Reproductive Health. 2020. PMID: 32138744.
\bibitem{ref21} Spencer CN. Health effects associated with exposure to intimate partner violence against women and childhood sexual abuse: a burden of proof study. Nature Medicine. 2023. PMID: 38081957.
\bibitem{ref22} Oram S. The Lancet Psychiatry Commission on intimate partner violence and mental health: advancing mental health services, research, and policy. The Lancet Psychiatry. 2022. PMID: 35569504.
\bibitem{ref23} Bacchus LJ. Interventions that prevent or respond to intimate partner violence against women and violence against children: a systematic review. The Lancet Public Health. 2024. PMID: 38702097.
\bibitem{ref24} Abramsky T. Findings from the SASA! Study: a cluster randomized controlled trial to assess the impact of a community mobilization intervention to prevent violence against women and reduce HIV risk in Kampala, Uganda. BMC medicine. 2014. PMID: 25248996.
\bibitem{ref25} Abramsky T. Ecological pathways to prevention: How does the SASA! community mobilisation model work to prevent physical intimate partner violence against women? BMC public health. 2016. PMID: 27084116.
\bibitem{ref26} Dunkle K. Effective prevention of intimate partner violence through couples training: a randomised controlled trial of Indashyikirwa in Rwanda. BMJ global health. 2020. PMID: 33355268.
\bibitem{ref27} Shaw B. Shifting Norms in Faith Communities to Reduce Intimate Partner Violence: Results from a Cluster Randomized Controlled Trial in Nigeria. Journal of interpersonal violence. 2023. PMID: 37329160.
\bibitem{ref28} Chatterji S. Community activism as a strategy to reduce intimate partner violence (IPV) in rural Rwanda: Results of a community randomised trial. Journal of global health. 2020. PMID: 32257154.
\bibitem{ref29} Jeong J. Effectiveness of the Moments that Matter parenting programme on caregiving practices, family violence, mental health and early child development in Kenya: a cluster-randomised controlled trial. BMJ global health. 2026. PMID: 42556995.
\bibitem{ref30} Ullman C. Interventions to prevent violence against women and girls globally: a global systematic review of reviews to update the RESPECT women framework. BMJ Public Health. 2025. PMID: 40017928.
\bibitem{ref31} Falb KL. Adapting, Sustaining and Scaling Violence against Women and Children Prevention Interventions in Low- and Middle-Income Countries: A Scoping Review on the Use of Implementation Science. Trauma, violence \& abuse. 2026. PMID: 40462272.
\bibitem{ref32} Doe Patience Fakornam. Couple-based preconception care as a gender-responsive strategy for improving health and well-being: a systematic review. Reproductive Health. 2026. PMID: 42316254.
\bibitem{ref33} Hawkins AJ. Exploring programmatic moderators of the effectiveness of marriage and relationship education programs: a meta-analytic study. Behavior Therapy. 2012. PMID: 22304880.
\bibitem{ref34} Bahrami-Samani S. Comparing the effects of premarital booklet- and video-based educations on the reproductive health literacy of engaged couples. Health Science Reports. 2024. PMID: 39377020.
\bibitem{ref35} Ashrafi E. Evaluating the Impact of Reproductive and Sexual Health Education on Knowledge and Attitudes of Couples in Premarital Counseling Classes: A Quasi-Experimental Study. Health Science Reports. 2026. PMID: 41773204.
\bibitem{ref36} Spencer CM. Online relationship education programs improve individual and relationship functioning: A meta-analytic review. Journal of Marital and Family Therapy. 2021. PMID: 33502054.
\bibitem{ref37} Robles TF. Marital quality and health: a meta-analytic review. Psychological Bulletin. 2014. PMID: 23527470.
\bibitem{ref38} O'Doherty L. Screening women for intimate partner violence in healthcare settings. Cochrane Database of Systematic Reviews. 2015. PMID: 26200817.
\bibitem{ref39} Feltner Cynthia. Screening for Intimate Partner Violence and for Caregiver Abuse of Older or Vulnerable Adults: An Evidence Report and Systematic Review for the US Preventive Services Task Force. JAMA. 2025. PMID: 40553460.
\bibitem{ref40} US Preventive Services Task Force. Screening for Intimate Partner Violence and Caregiver Abuse of Older or Vulnerable Adults: US Preventive Services Task Force Recommendation Statement. JAMA. 2025. PMID: 40553450.
\bibitem{ref41} Kalra Naira. Training healthcare providers to respond to intimate partner violence against women. Cochrane Database of Systematic Reviews. 2021. PMID: 34057734.
\bibitem{ref42} Betron. Effectiveness of a clinic-based counselling intervention on risk of experiencing intimate partner violence and reproductive coercion: a matched-pair cluster-controlled trial in Ebonyi and Sokoto, Nigeria. BMJ Global Health. 2025. PMID: 41151831.
\bibitem{ref43} World Health Organization. Responding to intimate partner violence and sexual violence against women: WHO clinical and policy guidelines. Geneva: World Health Organization; 2013. ISBN 978 92 4 154859 5. Available from: https://iris.who.int/handle/10665/85240; not PubMed-indexed.
\bibitem{ref44} World Health Organization. Health care for women subjected to intimate partner violence or sexual violence: a clinical handbook. Geneva: World Health Organization; 2014. WHO reference WHO/RHR/14.26. Available from: https://iris.who.int/handle/10665/136101; not PubMed-indexed.
\bibitem{ref45} Karakurt. Couples Therapy for Intimate Partner Violence: A Systematic Review and Meta-Analysis. Journal of marital and family therapy. 2016. PMID: 27377617.
\bibitem{ref46} McCollum. Couples treatment for interpersonal violence: a review of outcome research literature and current clinical practices. Violence and victims. 2008. PMID: 18624089.
\bibitem{ref47} Schneider. From scared to repaired: using an attachment-based perspective to understand situational couple violence. Journal of marital and family therapy. 2014. PMID: 25039387.
\bibitem{ref48} Chang. Health care interventions for intimate partner violence: what women want. Women's health issues. 2005. PMID: 15661584.
\bibitem{ref49} Robertson. Women and men's use of coercive control in intimate partner violence. Violence and victims. 2011. PMID: 21780535.
\bibitem{ref50} Swan. Factor structures for aggression and victimization among women who used aggression against male partners. Violence against women. 2012. PMID: 23012348.
\bibitem{ref51} Evans. "Even 'Daily' is Not Enough": How Well Do We Measure Domestic Violence and Abuse?-A Think-Aloud Study of a Commonly Used Self-Report Scale. Violence and victims. 2016. PMID: 26645540.
\bibitem{ref52} Skivington K. A new framework for developing and evaluating complex interventions: update of Medical Research Council guidance. BMJ. 2021. PMID: 34593508.
\bibitem{ref53} Moore G. Adapting interventions to new contexts-the ADAPT guidance. BMJ. 2021. PMID: 34344699.
\bibitem{ref54} Moore GF. Process evaluation of complex interventions: Medical Research Council guidance. BMJ. 2015. PMID: 25791983.
\bibitem{ref55} Horvat L. Cultural competence education for health professionals. Cochrane Database of Systematic Reviews. 2014. PMID: 24793445.
\bibitem{ref56} Vella Elizabeth. Does cultural competence training for health professionals impact culturally and linguistically diverse patient outcomes? A systematic review of the literature. Nurse Education Today. 2022. PMID: 35964378.
\bibitem{ref57} Hardy Billie-Jo. Systematic review of Indigenous cultural safety training interventions for healthcare professionals in Australia, Canada, New Zealand and the United States. BMJ Open. 2023. PMID: 37793931.
\bibitem{ref58} Tervalon M. Cultural humility versus cultural competence: a critical distinction in defining physician training outcomes in multicultural education. Journal of Health Care for the Poor and Underserved. 1998. PMID: 10073197.
\bibitem{ref59} Bresnahan Mary. Culturally safe healthcare: changing the lens from provider control to patient agency. Journal of Communication in Healthcare. 2024. PMID: 38426444.
\bibitem{ref60} Cano-Ibáñez Naomi. Physician-Patient Language Discordance and Poor Health Outcomes: A Systematic Scoping Review. Frontiers in Public Health. 2021. PMID: 33816420.
\bibitem{ref61} Chu Janet N. Association between language discordance and unplanned hospital readmissions or emergency department revisits: a systematic review and meta-analysis. BMJ Quality \& Safety. 2024. PMID: 38160059.
\bibitem{ref62} Heath Morten. Interpreter services and effect on healthcare - a systematic review of the impact of different types of interpreters on patient outcome. Journal of Migration and Health. 2023. PMID: 36816444.
\bibitem{ref63} Lauridsen Iben Gad. A systematic review of whether the number of linguistic errors in medical interpretation is associated with the use of professional vs ad hoc interpreters. Archives of Public Health. 2024. PMID: 39696623.
\bibitem{ref64} Reaume M. Patient-Physician Language Concordance and Cardiovascular Outcomes Among Patients With Hypertension. JAMA Network Open. 2025. PMID: 39969882.
\bibitem{ref65} Kosiyaporn Hathairat. Strengthening the migrant-friendliness of Thai health services through interpretation and cultural mediation: a system analysis. Global Health Research and Policy. 2020. PMID: 33372646.
\bibitem{ref66} König Andrea. A Systematic Scoping Review on Migrant Health Coverage in Thailand. Tropical Medicine and Infectious Disease. 2022. PMID: 36006258.
\bibitem{ref67} Phanwichatkul Titaree. The perceptions and practices of Thai health professionals providing maternity care for migrant Burmese women: An ethnographic study. Women and Birth. 2022. PMID: 34272187.
\bibitem{ref68} Buathong Napakkawat. Exploring symptoms perception and barriers to medication adherence among Thai Muslim patients with non-communicable diseases in a rural community in southern Thailand: a mixed-methods study. BMJ Open. 2024. PMID: 39653570.
\bibitem{ref69} Lohmae Userow. Muslim mothers' experiences in taking care of children with open heart surgery: A qualitative study in Southern Thailand. Belitung Nursing Journal. 2025. PMID: 40256382.
\bibitem{ref70} Pérez-Sánchez María. Access of migrant women to sexual and reproductive health services: A systematic review. Midwifery. 2024. PMID: 39243595.
\bibitem{ref71} Weschasat Tassanapan. The edutainment program on knowledge, perception, and uptake of cervical cancer screening among Muslim women in Southern Thailand: a quasi experimental study. BMC Public Health. 2024. PMID: 38971727.
\bibitem{ref72} Chuemchit Montakarn. Prevalence of Intimate Partner Violence in Thailand. Journal of Family Violence. 2018. PMID: 29904232.
\bibitem{ref73} Thananowan Nanthana. Intimate partner violence and factors associated with sexually transmitted infections among Thai women attending gynecology clinics. International Journal of STD \& AIDS. 2021. PMID: 33308089.
\bibitem{ref74} Thananowan Nanthana. Intimate partner violence among pregnant Thai women. Violence Against Women. 2008. PMID: 18408170.
\bibitem{ref75} Thitiyanviroj Benjaporn. Perceptions of Screening Women for Intimate Partner Violence Among Health Care Providers in Thailand. Nursing for Women's Health. 2024. PMID: 39536796.
\bibitem{ref76} Hayee Fusiyah. Sexual risk behaviors and influencing factors among Muslim adolescents on southern border of Thailand. International Journal of Adolescent Medicine and Health. 2020. PMID: 32549162.
\bibitem{ref77} Singkun Awirut. Risk of HIV/STIs among Muslim army conscripts in the three deep southern provinces of Thailand. Health Promotion Perspectives. 2021. PMID: 35079589.
\bibitem{ref78} Singkun Awirut. Sexual knowledge based on Islamic values and sexual risk behaviors of HIV/STIs among Thai Muslim army conscripts: A cross-sectional study. Belitung Nursing Journal. 2022. PMID: 37554485.
\bibitem{ref79} Ford Kathleen. Mental health in a conflict area: Migration, economic stress and religiosity in the three southernmost provinces of Thailand. International Journal of Social Psychiatry. 2017. PMID: 28024446.
\bibitem{ref80} Ford Kathleen. Mental health among youth in a conflict area: a cross-sectional study of the effects of parental absence, social media, family support and religious practice on psychiatric symptoms in Thailand's southernmost provinces. Conflict and Health. 2026. PMID: 41540416.
\bibitem{ref81} Vapattanawong Patama. A cross-sectional study on factors influencing adolescent depression in the civil-conflict areas of Thailand. BMC Public Health. 2026. PMID: 41872869.
\bibitem{ref82} Ford Kathleen. Coming of age in a conflict area: Mental health, education, employment, migration and family formation in the southernmost provinces of Thailand. International Journal of Social Psychiatry. 2018. PMID: 29417854.
\bibitem{ref83} Sirithammaphan Uraiwan. Barriers to measles mumps rubella vaccine acceptance in the three southern border provinces of Thailand. Clinical and Experimental Vaccine Research. 2023. PMID: 38025912.
\bibitem{ref84} Jinarong Taqwa. Muslim parents' beliefs and factors influencing complete immunization of children aged 0-5 years in a Thai rural community: a qualitative study. BMC Public Health. 2023. PMID: 37442996.
\bibitem{ref85} Napier AD. Culture and health. Lancet. 2014. PMID: 25443490.
\bibitem{ref86} Devakumar D. Racism, xenophobia, discrimination, and the determination of health. Lancet. 2022. PMID: 36502848.
\bibitem{ref87} Selvarajah S. Racism, xenophobia, and discrimination: mapping pathways to health outcomes. Lancet. 2022. PMID: 36502849.
\bibitem{ref88} Abubakar I. The UCL-Lancet Commission on Migration and Health: review of the state of progress. Lancet. 2026. PMID: 42314718.
\bibitem{ref89} Li S. Efficacy of culturally adapted interventions for common mental disorders in people of Chinese descent: a systematic review and meta-analysis. Lancet Psychiatry. 2023. PMID: 37208113.
\bibitem{ref90} Axinn WG. Community-Level Social Support Infrastructure and Adult Onset of Major Depressive Disorder in a South Asian Postconflict Setting. JAMA Psychiatry. 2022. PMID: 35080609.
\bibitem{ref91} Watt RG. Effectiveness and cost-effectiveness of a parenting programme to improve family wellbeing in England (TOGETHER): a multicentre, single-blind, randomised controlled trial. Lancet Public Health. 2026. PMID: 41887826.
\bibitem{ref92} Taft. Examining strength at home couples to prevent intimate partner violence on a military installation: A randomized controlled trial. Journal of consulting and clinical psychology. 2024. PMID: 38206858.
\bibitem{ref93} Ghuman SJ. Women's autonomy and child survival: a comparison of Muslims and non-Muslims in four Asian countries. Demography. 2003. PMID: 12962056.
\bibitem{ref94} Chen. Intimate partner violence: counseling, community resources, and legal issues for IPV victims and perpetrators. FP essentials. 2013. PMID: 24053261.
\bibitem{ref95} Decker. Implementing Trauma-Informed Partner Violence Assessment in Family Planning Clinics. Journal of women's health. 2017. PMID: 28375750.
\bibitem{ref96} MacMillan. Approaches to screening for intimate partner violence in health care settings: a randomized trial. JAMA. 2006. PMID: 16882959.
\bibitem{ref97} Ford-Gilboe. Development of a brief measure of intimate partner violence experiences: the Composite Abuse Scale (Revised)-Short Form (CASR-SF). BMJ open. 2016. PMID: 27927659.
\bibitem{ref98} Hegarty. Validation of the Coercive-Composite Abuse Scale (C-CAS) in an Australian sample of women experiencing intimate partner abuse. BMJ open. 2026. PMID: 42323273.
\bibitem{ref99} Wilson. Psychometric Properties of the Coercion in Intimate Partner Relationships Scale. Assessment. 2023. PMID: 34167331.
\bibitem{ref100} Sarac A. Cultural adaptation and pilot evaluation of the Couple Learning Program (CLP) among young Turkish couples: A mixed-method study. Acta Psychologica. 2026. PMID: 42442188.
\bibitem{ref101} Welland C. Culturally specific treatment for partner-abusive Latino men: a qualitative study to identify and implement program components. Violence and Victims. 2010. PMID: 21287968.
\bibitem{ref102} Sikoki B. A qualitative study on perceptions of adolescents' sexual and reproductive health education in Yogyakarta, Indonesia. International Journal of Adolescent Medicine and Health. 2024. PMID: 39395200.
\bibitem{ref103} Fauk NK. Cultural and religious determinants of HIV transmission: A qualitative study with people living with HIV in Belu and Yogyakarta, Indonesia. PLoS ONE. 2021. PMID: 34780506.
\bibitem{ref104} Grisurapong S. Establishing a one-stop crisis center for women suffering violence in Khonkaen hospital, Thailand. International Journal of Gynaecology and Obstetrics. 2002. PMID: 12429436.
\bibitem{ref105} Nagae M. [Domestic violence as a women's health issue--activities of health professionals in Thailand]. Japanese Journal of Public Health (Nihon Koshu Eisei Zasshi). 2004. PMID: 15162975.
\bibitem{ref106} Williamson HC. Does premarital education decrease or increase couples' later help-seeking? Journal of Family Psychology. 2014. PMID: 24294929.
\bibitem{ref107} Williamson HC. Barriers and facilitators of relationship help-seeking among low-income couples. Journal of Family Psychology. 2019. PMID: 30489129.
\bibitem{ref108} Oduenyi. Integrating gender-based violence first-line response into family planning and antenatal care services: an assessment of feasibility, appropriateness and acceptability to healthcare providers in Sokoto and Ebonyi States, Nigeria. BMC Health Services Research. 2025. PMID: 41267027.
\bibitem{ref109} Rhodes. The anatomy of a community health center system-level intervention for intimate partner violence. Journal of Urban Health. 2014. PMID: 23917943.
\bibitem{ref110} Miller Elizabeth. A family planning clinic-based intervention to address reproductive coercion: a cluster randomized controlled trial. Contraception. 2016. PMID: 26892333.
\bibitem{ref111} Zaher. Effect of domestic violence training: systematic review of randomized controlled trials. Canadian Family Physician. 2014. PMID: 25022633.
\bibitem{ref112} Jacubowski N. Marriage is not a safe place: heterosexual marriage and HIV-related vulnerability in Indonesia. Culture, Health \& Sexuality. 2008. PMID: 18038283.
\bibitem{ref113} Michau L. SASA! Together: An evolution of the SASA! approach to prevent violence against women. Evaluation and program planning. 2021. PMID: 33578229.
\bibitem{ref114} Stern E. Lessons learned from implementing Indashyikirwa in Rwanda - an adaptation of the SASA! approach to prevent and respond to intimate partner violence. Evaluation and program planning. 2018. PMID: 30125773.
\bibitem{ref115} Clyde TL. Revising Premarital Relationship Interventions for the Next Generation. Journal of Marital and Family Therapy. 2020. PMID: 30725473.
\bibitem{ref116} Silliman B. Improving practice in marriage preparation. Journal of Sex \& Marital Therapy. 1999. PMID: 10081741.
\bibitem{ref117} Paiboonsukwong Kittiphong. Thalassemia in Thailand. Hemoglobin. 2022. PMID: 35950590.
\bibitem{ref118} Zust BL. 10-Year Study of Christian Church Support for Domestic Violence Victims: 2005-2015. Journal of Interpersonal Violence. 2021. PMID: 33729071.
\bibitem{ref119} Hegarty K et al. Protocol for a randomised controlled trial of a web-based healthy relationship tool and safety decision aid for women experiencing domestic violence (I-DECIDE). BMC Public Health. 2015; PMID: 26231225.
\bibitem{ref120} Hegarty K et al. An online healthy relationship tool and safety decision aid for women experiencing intimate partner violence (I-DECIDE): a randomised controlled trial. The Lancet Public Health. 2019; PMID: 31155223.
\bibitem{ref121} Tarzia L et al. I-DECIDE: An Online Intervention Drawing on the Psychosocial Readiness Model for Women Experiencing Domestic Violence. Women's Health Issues. 2016; PMID: 26362841.
\bibitem{ref122} Tarzia L et al. Methodological and Ethical Challenges in a Web-Based Randomized Controlled Trial of a Domestic Violence Intervention. Journal of Medical Internet Research. 2017; PMID: 28351830.
\bibitem{ref123} Glass N et al. A safety app to respond to dating violence for college women and their friends: the MyPlan study randomized controlled trial protocol. BMC Public Health. 2015; PMID: 26350482.
\bibitem{ref124} Glass NE et al. Longitudinal Impact of the myPlan App on Health and Safety Among College Women Experiencing Partner Violence. Journal of Interpersonal Violence. 2022; PMID: 33576291.
\bibitem{ref125} Bloom TL et al. Concerned friends of intimate partner violence survivors: results from the myPlan randomized controlled trial on college campuses. BMC Public Health. 2023; PMID: 37259087.
\bibitem{ref126} Glass N et al. Effectiveness of the myPlan Teen App, a Digital Healthy Relationship and Safety Planning Intervention With Adolescent Aged 15-17 Years. Journal of Adolescent Health. 2024; PMID: 39066749.
\bibitem{ref127} Decker MR et al. Adapting the myPlan safety app to respond to intimate partner violence for women in low and middle income country settings: app tailoring and randomized controlled trial protocol. BMC Public Health. 2020; PMID: 32471469.
\bibitem{ref128} Decker MR et al. Safety decision-making and planning mobile app for intimate partner violence prevention and response: randomised controlled trial in Kenya. BMJ Global Health. 2020; PMID: 32675229.
\bibitem{ref129} Ford-Gilboe M et al. A tailored online safety and health intervention for women experiencing intimate partner violence: the iCAN Plan 4 Safety randomized controlled trial protocol. BMC Public Health. 2017; PMID: 28327116.
\bibitem{ref130} Ford-Gilboe M et al. Longitudinal impacts of an online safety and health intervention for women experiencing intimate partner violence: randomized controlled trial. BMC Public Health. 2020; PMID: 32098633.
\bibitem{ref131} O'Campo P et al. Design and Development of a Suite of Intimate Partner Violence Screening and Safety Planning Web Apps: User-Centered Approach. Journal of Medical Internet Research. 2021; PMID: 34931998.
\bibitem{ref132} Chuemchit M et al. Development and Preliminary Evaluation of iCanPlan: A Mobile Health Application for Intimate Partner Violence Prevention in Thailand. International Journal of Environmental Research and Public Health. 2026; PMID: 42196762.
\bibitem{ref133} Berry-Cabán CS et al. Empowering Safety: The Case for a Decision-making App for Military Partners Facing Intimate Partner Violence. Military Medicine. 2026; PMID: 42560214.
\bibitem{ref134} Ahmad F et al. Computer-assisted screening for intimate partner violence and control: a randomized trial. Annals of Internal Medicine. 2009; PMID: 19487706.
\bibitem{ref135} Lenert L et al. Electronic Health Record-Based Screening for Intimate Partner Violence: A Cluster Randomized Clinical Trial. JAMA Network Open. 2024; PMID: 39088215.
\bibitem{ref136} Wood SN et al. Risk for Severe Intimate Partner Violence in Nairobi's Informal Settlements: Tailoring the Danger Assessment to Kenya. Global Health, Science and Practice. 2024; PMID: 38290753.
\bibitem{ref137} El Morr C et al. Effectiveness of ICT-based intimate partner violence interventions: a systematic review. BMC Public Health. 2020; PMID: 32894115.
\bibitem{ref138} Anderson EJ et al. Web-Based and mHealth Interventions for Intimate Partner Violence Victimization Prevention: A Systematic Review. Trauma, Violence, \& Abuse. 2021; PMID: 31742475.
\bibitem{ref139} Ali K et al. Online Peer-to-Peer Support for Young People With Mental Health Problems: A Systematic Review. JMIR Mental Health. 2015; PMID: 26543923.
\bibitem{ref140} Marshall P et al. Understanding the Impacts of Online Mental Health Peer Support Forums: Realist Synthesis. JMIR Mental Health. 2024; PMID: 38722680.
\bibitem{ref141} Marshall P et al. Understanding Safety in Online Mental Health Forums: Realist Evaluation. JMIR Mental Health. 2025; PMID: 40577699.
\bibitem{ref142} Perry A et al. Suicidal behaviours and moderator support in online health communities: a scoping review. BMJ Open. 2021; PMID: 34193497.
\bibitem{ref143} Perowne R et al. Barriers and enablers to the moderation of self-harm content for a young person's online forum. Journal of Mental Health. 2024; PMID: 35574666.
\bibitem{ref144} Deng D et al. The Role of Moderators in Facilitating and Encouraging Peer-to-Peer Support in an Online Mental Health Community: A Qualitative Exploratory Study. Journal of Technology in Behavioral Science. 2023; PMID: 36810998.
\bibitem{ref145} Lerman K et al. Safe spaces or toxic places? Content moderation and social dynamics of online eating disorder communities. EPJ Data Science. 2025; PMID: 40726606.
\bibitem{ref146} Bender JL et al. How Cancer Online Support Groups Work, for Whom, and in What Circumstances: Realist Review. Journal of Medical Internet Research. 2026; PMID: 42127226.
\bibitem{ref147} Chu TH et al. Online Social Support for Intimate Partner Violence Victims in China: Quantitative and Automatic Content Analysis. Violence Against Women. 2021; PMID: 32339059.
\bibitem{ref148} Rempel E et al. Intimate partner violence: a review of online interventions. Informatics for Health \& Social Care. 2019; PMID: 29537928.
\bibitem{ref149} Rogers MM et al. Technology-Facilitated Abuse in Intimate Relationships: A Scoping Review. Trauma, Violence \& Abuse. 2023; PMID: 35537445.
\bibitem{ref150} Henry N et al. Technology-Facilitated Domestic and Sexual Violence: A Review. Violence Against Women. 2020; PMID: 32998673.
\bibitem{ref151} Henry N et al. Technology-Facilitated Domestic Violence Against Immigrant and Refugee Women: A Qualitative Study. Journal of Interpersonal Violence. 2022; PMID: 33719681.
\bibitem{ref152} Weekes CJ et al. Cyberstalking Perpetrators and Their Methods: A Systematic Literature Review. Trauma, Violence \& Abuse. 2026; PMID: 40271872.
\bibitem{ref153} Irving M et al. Post-separation Parenting Apps as Tools for Control and Resistance. Violence Against Women. 2026; PMID: 41823642.
\bibitem{ref154} Schlosser AE et al. Self-disclosure versus self-presentation on social media. Current Opinion in Psychology. 2020; PMID: 31357096.
\bibitem{ref156} Pretorius C et al. Young People's Online Help-Seeking and Mental Health Difficulties: Systematic Narrative Review. Journal of Medical Internet Research. 2019; PMID: 31742562.
\bibitem{ref157} Douglass CH et al. Social Media and Online Digital Technology Use Among Muslim Young People and Parents: Qualitative Focus Group Study. JMIR Pediatrics and Parenting. 2022; PMID: 35536616.
\bibitem{ref158} Robbins KC et al. The Independent Living Donor Advocate: An Essential Role for Living Kidney Donation. Nephrology Nursing Journal. 2014; PMID: 26287054.
\bibitem{ref159} Rudow DL et al. The Psychosocial and Independent Living Donor Advocate Evaluation and Post-surgery Care of Living Donors. Journal of Clinical Psychology in Medical Settings. 2015; PMID: 26293351.
\bibitem{ref160} Schenck D et al. Ethical Analysis and Policy Recommendations Regarding Domino Liver Transplantation. Transplantation. 2018; PMID: 29708521.
\bibitem{ref161} Sethi N et al. Ethics and the facial plastic surgeon. European Archives of Oto-Rhino-Laryngology. 2016; PMID: 26254909.
\bibitem{ref162} Pope KS et al. Dual relationships in psychotherapy. Ethics \& Behavior. 1991; PMID: 11649348.
\bibitem{ref163} Scopelliti J et al. Dual relationships in mental health practice: issues for clinicians in rural settings. Australian and New Zealand Journal of Psychiatry. 2004; PMID: 15555031.
\bibitem{ref164} Messing JT et al. The average predictive validity of intimate partner violence risk assessment instruments. Journal of Interpersonal Violence. 2013; PMID: 23262817.
\bibitem{ref165} Wang PL et al. Assessing the Danger: Validation of Taiwan Intimate Partner Violence Danger Assessment. Journal of Interpersonal Violence. 2015; PMID: 25315482.
\bibitem{ref166} Messing JT et al. Validation and adaptation of the danger assessment-5: A brief intimate partner violence risk assessment. Journal of Advanced Nursing. 2017; PMID: 28921610.
\bibitem{ref167} Messing JT et al. Accounting for Multiple Nonfatal Strangulation in Intimate Partner Violence Risk Assessment. Journal of Interpersonal Violence. 2022; PMID: 33280504.
\bibitem{ref168} Gerth J et al. [The Ontario Domestic Assault Risk Assessment (ODARA) - Validity und authorised German translation of an intimate partner violence screening tool]. Fortschritte der Neurologie-Psychiatrie. 2014; PMID: 25383928.
\bibitem{ref169} Jolliffe Simpson AD et al. Unpacking Multiagency Structured Professional Judgment Risk Assessments for Family Violence. Journal of Interpersonal Violence. 2023; PMID: 36710516.
\bibitem{ref170} Belfrage H et al. Assessment and management of risk for intimate partner violence by police officers using the Spousal Assault Risk Assessment Guide. Law and Human Behavior. 2012; PMID: 22471386.
\bibitem{ref171} Hogan NR et al. Investigating racial disparities in violence risk assessment using the Spousal Assault Risk Assessment Guide-Version 3. Psychological Assessment. 2024; PMID: 38512165.
\bibitem{ref172} Sanderson D et al. "Increasing Warm Handoffs: Optimizing Community Based Referrals in Primary Care Using QI Methodology. Journal of Primary Care \& Community Health. 2021; PMID: 34109884.
\bibitem{ref173} Loo S et al. Evaluating a social risk screening and referral program in an urban safety-net hospital emergency department. Journal of the American College of Emergency Physicians Open. 2023; PMID: 36704207.
\bibitem{ref174} Stevens GJ et al. Efficacy of a short message service brief contact intervention (SMS-SOS) in reducing repetition of hospital-treated self-harm. British Journal of Psychiatry. 2024; PMID: 38083861.
\bibitem{ref175} Radin AK et al. Comparative effectiveness of safety planning intervention with instrumental support calls (ISC) versus caring contacts (CC) two-way text messages: the SPARC Trial protocol. Contemporary Clinical Trials. 2023; PMID: 37321352.
\bibitem{ref176} Radin AK et al. Comparative effectiveness of two versions of a caring contacts intervention in healthcare providers, staff and patients. Journal of Affective Disorders. 2023; PMID: 36963515.
\bibitem{ref179} Chambers RA et al. Implementation Fidelity and Theory-Informed Dose Effects of a Teen Pregnancy Prevention Program for Native American Youth. Prevention Science. 2023; PMID: 37191932.
\bibitem{ref182} Roozenbeek J et al. Psychological inoculation improves resilience against misinformation on social media. Science advances. 2022; PMID: 36001675.
\bibitem{ref183} Appel RE et al. Psychological inoculation improves resilience to and reduces willingness to share vaccine misinformation. Scientific reports. 2025; PMID: 40820144.
\bibitem{ref184} Biddlestone M et al. Video inoculation against election misinformation across 12 EU nations. Communications psychology. 2026; PMID: 41845063.
\bibitem{ref185} Basol M et al. Good News about Bad News: Gamified Inoculation Boosts Confidence and Cognitive Immunity Against Fake News. Journal of cognition. 2020; PMID: 31934684.
\bibitem{ref186} Roozenbeek J et al. Technique-based inoculation against real-world misinformation. Royal Society open science. 2022; PMID: 35600423.
\bibitem{ref187} Roozenbeek J et al. Disentangling Item and Testing Effects in Inoculation Research on Online Misinformation: Solomon Revisited. Educational and psychological measurement. 2021; PMID: 37929261.
\bibitem{ref188} Kang Y et al. Educational and skills-based interventions for preventing relationship and dating violence in adolescents and young adults. The Cochrane database of systematic reviews. 2026; PMID: 42535453.
\bibitem{ref189} Kettrey HH et al. Can Attitudes Serve as Proxies for Behavioral Outcomes of Dating Violence Prevention Programs? Broader Lessons From a Pilot Evaluation of the Relationship Education Project. Violence and victims. 2023; PMID: 37011949.
\bibitem{ref190} Miller E et al. One-year follow-up of a coach-delivered dating violence prevention program: a cluster randomized controlled trial. American journal of preventive medicine. 2013; PMID: 23790995.
\bibitem{ref191} Abebe KZ et al. A cluster-randomized trial of a middle school gender violence prevention program: Design, rationale, and sample characteristics. Contemporary clinical trials. 2017; PMID: 28821469.
\bibitem{ref192} Cherrier C et al. [Supporting the applicability and transferability of the "Sortir Ensemble \& Se Respecter" program in France]. Sante publique (Vandoeuvre-les-Nancy, France). 2024; PMID: 38580464.
\bibitem{ref193} Bannon RS et al. The Bystander Approach to Sexual Assault Risk Reduction: Effects on Risk Recognition, Perceived Self-Efficacy, and Protective Behavior. Violence and victims. 2017; PMID: 28234197.
\bibitem{ref194} Peeren S et al. "Counteract the gaslighting" - a thematic analysis of open-ended responses about what women survivors of intimate partner sexual violence need from services. BMC women's health. 2024; PMID: 38336660.
\bibitem{ref195} Barnes M et al. Being silenced, loneliness and being heard: understanding pathways to intimate partner violence \& abuse in young adults. a mixed-methods study. BMC public health. 2022; PMID: 35974354.
\bibitem{ref196} Whitton SW et al. Disclosure and Help-Seeking Experiences of Sexual and Gender Minority Victims of Intimate Partner Violence: A Mixed-Methods Study. Journal of interpersonal violence. 2024; PMID: 37882155.
\bibitem{ref197} Brunton R et al. Further psychometric examination of the Intimate Partner Violence Internalized Stigma Scale (IPVIS; Brunton \& Harris, 2023): Examining criterion-related validity. Acta psychologica. 2024; PMID: 39536684.
\bibitem{ref198} Tian CY et al. Generic Health Literacy Measurements for Adults: A Scoping Review. International journal of environmental research and public health. 2020; PMID: 33114157.
\bibitem{ref199} Senahad N et al. Development of health literacy assessment tool for 9-10 Years old children in Thailand. Public health in practice (Oxford, England). 2023; PMID: 37727706.
\bibitem{ref200} Couture-Carron A et al. One Size Doesn't Fit All: Dating Abuse Against Women From the Perspective of South Asian Muslim Youth in Canada. Journal of interpersonal violence. 2017; PMID: 26289454.
\bibitem{ref201} Thongpriwan V et al. Reflections on attitudes, experiences, and vulnerability of intimate partner violence among Southeast Asian college women living in United States. Asian journal of psychiatry. 2015; PMID: 26442989.
\bibitem{ref202} Bell CA et al. Humility, Forgiveness, and Emerging Adult Female Romantic Relationships. Journal of marital and family therapy. 2019; PMID: 29073326.
\bibitem{ref203} Koetke J et al. Intellectual humility is reliably associated with constructive responses to conflict. PloS one. 2024; PMID: 39240981.
\bibitem{ref204} Lehmann M et al. Intellectual Humility Predicts Empathic Accuracy and Empathic Resilience. Personality \& social psychology bulletin. 2026; PMID: 39902736.
\bibitem{ref205} Jongman-Sereno KP et al. Intellectual Humility and Responsiveness to Public Health Recommendations. Personality and individual differences. 2023; PMID: 37426514.
\bibitem{ref206} Foronda C et al. A Theory of Cultural Humility. Journal of transcultural nursing. 2020; PMID: 31516091.
\bibitem{ref207} Lekas HM et al. Rethinking Cultural Competence: Shifting to Cultural Humility. Health services insights. 2020; PMID: 33424230.
\bibitem{ref208} Sotto-Santiago S et al. S.E.L.F.: Lessons in institutional cultural humility approach from an academic health center. Dialogues in health. 2026; PMID: 42656867.
\bibitem{ref209} MacBeth A et al. Exploring compassion: a meta-analysis of the association between self-compassion and psychopathology. Clinical psychology review. 2012; PMID: 22796446.
\bibitem{ref210} Millard LA et al. The effectiveness of compassion focused therapy with clinical populations: A systematic review and meta-analysis. Journal of affective disorders. 2023; PMID: 36649790.
\bibitem{ref211} Baker LR et al. Self-compassion and relationship maintenance: the moderating roles of conscientiousness and gender. Journal of personality and social psychology. 2011; PMID: 21280964.
\bibitem{ref212} Kaya F et al. The predictive effect of self-compassion on relationship satisfaction and conflict resolution styles in romantic relationships in nursing students. Nursing forum. 2022; PMID: 35245387.
\bibitem{ref213} Williams Z et al. Relationship satisfaction in Black heterosexual couples: The role of self-compassion and openness. Journal of marital and family therapy. 2023; PMID: 37752743.
\bibitem{ref214} Godfrey DA et al. Empathy Mediates the Relations between Working Memory and Perpetration of Intimate Partner Violence and Aggression. Behavioral sciences (Basel, Switzerland). 2020; PMID: 32150915.
\bibitem{ref215} Brassard A et al. Childhood Sexual Abuse, Dyadic Empathy, and Intimate Partner Violence Among Men Seeking Psychological Help. Journal of interpersonal violence. 2022; PMID: 35089108.
\bibitem{ref216} Buitelaar NJL et al. ADHD in Childhood and/or Adulthood as a Risk Factor for Domestic Violence or Intimate Partner Violence: A Systematic Review. Journal of attention disorders. 2020; PMID: 25995243.
\bibitem{ref217} Boonyananth N et al. A Systematic Review of the Psychosocial Mechanism Underlying the Pathway from Intra-Familial Victimization in Childhood to Intimate Relationship Violence in Adolescence and Adulthood. Trauma, violence \& abuse. 2026; PMID: 40035136.
\bibitem{ref218} Gopalakrishnan L et al. Psychometric Validation of Trust, Commitment, and Satisfaction Scales to Measure Marital Relationship Quality Among Newly Married Women in Nepal. International journal of environmental research and public health. 2025; PMID: 41007600.
\bibitem{ref219} Körün AB et al. Even Though the Long Distance: Are We Still Going on? Dyadic Trust, Relationship Maintenance Behaviors, and Relationship Quality Among Emerging Adulthoods. Scandinavian journal of psychology. 2026; PMID: 41235925.
\bibitem{ref220} Gonçalves JPB et al. The role of religiosity and spirituality in interpersonal violence: a systematic review and meta-analysis. Revista brasileira de psiquiatria (Sao Paulo, Brazil). 2023; PMID: 36331229.
\bibitem{ref221} Simonič B et al. The Power of Women's Faith in Coping with Intimate Partner Violence: Systematic Literature Review. Journal of religion and health. 2021; PMID: 33704630.
\bibitem{ref223} Ripley JS et al. Obtaining religious and spiritual competencies for relationship therapy: Outcomes of Competency Addressing Religion and Spirituality (CARS) training. Psychotherapy (Chicago, Ill.). 2026; PMID: 40705620.
\bibitem{ref224} Aziz AFA et al. Knowledge improvement on premarital HIV counseling with PAUSE flipchart versus conventional method: a cluster randomised controlled trial. BMC primary care. 2026; PMID: 42420861.
\bibitem{ref225} Mustafa Y et al. Islam and the four principles of medical ethics. Journal of Medical Ethics. 2014; PMID: 23975951.
\bibitem{ref226} Padela AI et al. Muslim patients and cross-gender interactions in medicine: an Islamic bioethical perspective. Journal of Medical Ethics. 2011; PMID: 21041237.
\bibitem{ref227} Doedes E et al. Moral universe of Muslim healthcare practitioners in the UK: balancing Islamic and secular ethics in palliative and end-of-life care. Journal of Medical Ethics. 2026; PMID: 41651650.
\bibitem{ref228} Padela AI et al. Islamic medical ethics: a primer. Bioethics. 2007; PMID: 17845488.
\bibitem{ref229} Alsolamy S et al. Islamic views on artificial nutrition and hydration in terminally ill patients. Bioethics. 2014; PMID: 22845721.
\bibitem{ref230} Daar AS et al. Bioethics for clinicians: 21. Islamic bioethics. CMAJ (Canadian Medical Association Journal). 2001; PMID: 11202669.
\bibitem{ref231} Bhowmik B et al. Faith-Based Lifestyle Intervention for Diabetes Prevention Among Adults in Bangladesh: A Cluster Randomized Clinical Trial. JAMA Network Open. 2025; PMID: 41114978.
\bibitem{ref232} Zoellner LA et al. Lay-Led Intervention for War and Refugee Trauma: A Randomized Clinical Trial. JAMA Network Open. 2024; PMID: 39186273. Two errata published: JAMA Netw Open. 2024;7(9):e2440198 and 2024;7(10):e2446257.
\bibitem{ref233} Gendler Y et al. Exploring Cultural and Religious Effects on HPV Vaccination Decision Making Using a Web-Based Decision Aid: A Quasi-experimental Study. Medical Decision Making. 2024; PMID: 38600776.
\bibitem{ref234} Al-Sayaghi KM et al. Stressors, Needs and Satisfaction of Families of Critically Ill Arabic Patients: A Systematic Review. Nursing in Critical Care. 2026; PMID: 41379069.
\bibitem{ref235} Mohamed Mahdi Z et al. Gender-Based Violence Against Migrant Women From Islamic Background: A Systematic Review of International Studies. Trauma, Violence \& Abuse. 2025; PMID: 41139992.
\bibitem{ref236} Chamsi-Pasha H et al. Assisted reproductive technology: Islamic Sunni perspective. Human Fertility. 2015; PMID: 25660098.
\bibitem{ref237} Pan SW et al. Health Effects of Religion, Spirituality, and Supernatural Beliefs in Mainland China: A Systematic Review. Journal of Religion and Health. 2022; PMID: 35347576.
\bibitem{ref238} Abdel-Khalek AM et al. The Relationship Between Religiosity and Anxiety: A Meta-analysis. Journal of Religion and Health. 2019; PMID: 31309442.
\bibitem{ref239} Cleary M et al. Drivers and consequences of self-immolation in parts of Iran, Iraq and Uzbekistan: A systematic review of qualitative evidence. Burns. 2021; PMID: 31928787.
\bibitem{ref240} Mittal SO et al. Faith, fasting, and well-being: Emirates Neurology Society consensus guidelines on safe Ramadan fasting in Parkinson's disease. Frontiers in Neurology. 2025; PMID: 41458120.
\bibitem{ref241} Packer S et al. Informed consent with a focus on islamic views. The Journal of IMA (Islamic Medical Association of North America) (professional-association journal). 2011; PMID: 23610513.
\bibitem{ref242} Arawi TA et al. The muslim physician and the ethics of medicine. The Journal of IMA (Islamic Medical Association of North America) (professional-association journal). 2010; PMID: 23864762.
\bibitem{ref243} Irfan B et al. Considering Islamic Frameworks to Infectious Disease Prevention. Open Forum Infectious Diseases. 2025; PMID: 41141460.
\end{thebibliography}
\end{document}
